\documentclass[11pt,letterpaper]{article}

\usepackage[margin=1in]{geometry}
\usepackage[T1]{fontenc}
\usepackage{lmodern}
\usepackage{microtype}
\usepackage{amsmath,amssymb,amsthm}
\usepackage{bbm}
\usepackage{bm}
\usepackage{graphicx}
\usepackage[x11names]{xcolor}
\usepackage{xspace}
\usepackage[ruled]{algorithm2e}
\usepackage[normalem]{ulem}
\usepackage[authoryear,round]{natbib}
\usepackage{url}
\usepackage{hyperref}
\hypersetup{
  colorlinks=true,
  linkcolor=blue,
  citecolor=blue,
  urlcolor=magenta,
  pdftitle={Towards a Theory of Learnability via Debate},
  pdfauthor={Tosca Lechner, Vinayak Pathak, Sivan Sabato, Ruth Urner}
}
\bibliographystyle{plainnat}
\bibpunct{(}{)}{;}{a}{,}{,}

\newtheorem{theorem}{Theorem}
\newtheorem{lemma}[theorem]{Lemma}
\theoremstyle{definition}
\newtheorem{definition}[theorem]{Definition}
\newtheorem{example}{Example}
\theoremstyle{plain}
\newtheorem{claim}[theorem]{Claim}

% caligraphic letters
  \newcommand{\Acal}{{\mathcal A}}
  \newcommand{\A}{{\mathcal A}}
  \newcommand{\Bcal}{{\mathcal B}}
  \newcommand{\B}{{\mathcal B}}
  \newcommand{\Ccal}{{\mathcal C}}
  \newcommand{\C}{{\Ccal}}
  \newcommand{\Dcal}{{\mathcal D}}
  \newcommand{\Fcal}{{\mathcal F}}
  \newcommand{\F}{{\mathcal F}}
  \newcommand{\Hcal}{{\mathcal H}}
  \renewcommand{\H}{{\Hcal}}
  \newcommand{\Ical}{{\mathcal I}}
  \newcommand{\Pcal}{{\mathcal P}}
  \renewcommand{\P}{{\mathcal P}}
  \newcommand{\Qcal}{{\mathcal Q}}
  \newcommand{\Ocal}{{\mathcal O}}
  \newcommand{\Ncal}{{\mathcal N}}
  \newcommand{\Tcal}{{\mathcal T}}
  \newcommand{\Ucal}{{\mathcal U}}
  \newcommand{\Vcal}{{\mathcal V}}
  \newcommand{\V}{{\mathcal V}}
  \newcommand{\Scal}{{\mathcal S}}
  \newcommand{\Wcal}{{\mathcal W}}
  \newcommand{\Xcal}{{\mathcal X}}
  \newcommand{\X}{{\mathcal X}}
  \newcommand{\Z}{{\mathcal Z}}
  \newcommand{\Y}{{\mathcal Y}}
  \newcommand{\Ycal}{{\mathcal Y}}
  \newcommand{\Zcal}{{\mathcal Z}}
  \newcommand{\ol}{\overline}
  \newcommand{\zero}{\bf{0}}

%Losses and mistake counts
  \newcommand{\Lo}{{\mathcal L}}
  \newcommand{\tLo}{{{\mathcal L}_P}}
  \newcommand{\fLo}{{{\mathcal L}_R}}
  \renewcommand{\Re}{{\mathcal R}}
  \newcommand{\tRe}{{{\mathcal R}_P}}
  \newcommand{\fRe}{{{\mathcal R}_R}}
  \newcommand{\Mi}{{\mathcal M}}
  \newcommand{\mi}{\mathrm{m}}


 \newcommand{\ie}{\textsl{i.e.\ }} % id est
 \newcommand{\iid}{\textsl{i.i.d.\ }} % id est
 \newcommand{\eg}{e.g.\ }          % e.g.
 \newcommand{\ex}{\mathbb{E}}
 \newcommand{\reals}{\mathbb{R}}
 \newcommand{\naturals}{\mathbb{N}}

\newcommand{\indct}{\mathbbm{1}}


 %certificates and debaters
  \newcommand{\dt}{\texttt{Proposer}\xspace}
  \newcommand{\df}{\texttt{Responder}\xspace}

  \renewcommand{\c}{c}
  \newcommand{\ct}{{p}}
  \newcommand{\cf}{{r}}
  \renewcommand{\sc}{\bar{c}}
  \newcommand{\hy}{{\hat{y}}}
  \newcommand{\vy}{{\widetilde{y}}}
  \newcommand{\ty}{y}

  %trees
  \newcommand{\T}{{\mathcal T}}
  \newcommand{\depth}{\mathrm{depth}}
  \newcommand{\dldim}{\mathrm{DLdim}}
 \newcommand{\debatedimname}{debate learning dimension\xspace}


% comments
\newcommand{\todo}[1]{{\color{red}todo: #1}}
\newcommand{\ruth}[1]{{\color{purple}Ruth: #1}}
\newcommand{\tosca}[1]{{\color{teal}Tosca: #1}}
\newcommand{\sivan}[1]{{\color{orange}Sivan: #1}}
\newcommand{\vinayak}[1]{{\color{olive}Vinayak: #1}}

\newcommand{\ctd}{\mathrm{\ct dim}}
\newcommand{\cfd}{\mathrm{\cf dim}}


\DeclareMathOperator*{\argmin}{arg\,min}
\DeclareMathOperator*{\argmax}{arg\,max}

\title{Towards a Theory of Learnability via Debate}
% Affiliations will be added later.
\author{
  Tosca Lechner \and
  Vinayak Pathak \and
  Sivan Sabato \and
  Ruth Urner
}
\date{}
\pagestyle{plain}

\begin{document}

\maketitle

\begin{abstract}%
  Debate has been recently proposed as a training protocol that is hoped to
produce more aligned models even as model capabilities scale. The idea is
that the human, instead of providing supervision in the form of labels,
judges a debate between two models that compete to prove the correctness of
opposite labels. It can be shown that functions of a much higher computational complexity  can be
correctly judged in this way, compared to the complexity of functions that can be directly evaluated using a computationally weak evaluator. This suggests that perhaps debate can be
used to teach very complicated functions to models. However, previous work
has largely ignored the fact that this is a \emph{training} protocol, and
simply analyzes the game being played between the two debaters for a known
judge. We take the first steps towards understanding the learning problem
implicit within the protocol, by formulating it as an online learning
problem. We show that once we formulate the learning problem precisely,
there are cases where the best way to learn given access to a judge is
different from running an adversarial debate. This motivates us to define a
cooperative version of the protocol, whose learnability we characterize
using a combinatorial dimension. We also consider the possibility of the
human judge being fallible and occasionally judging the debate incorrectly.
We show that a small number of errors in judgement can lead to a large error
in learning.
%
\end{abstract}

\noindent\textbf{Keywords.} Learning via debate, online learning, interactive learning

\section{Introduction}
As AI models grow in capabilities, how do we ensure that they continue to act in the interest of humans? A long line of work~\citep{yudkowsky2008artificial,soares2015corrigibility,russell2015research,amodei2016concrete,everitt2017corrupted,soares2017agent,armstrong2018occam,soares2018value,hubinger2019risks,taylor2020alignment,skalse2022reward} has argued that this is a hard problem, and recent events such as the infamous HuggingFace incident~\citep{openai2026huggingface} have demonstrated the difficulties in practice. 

One way to study this problem more formally has been termed \emph{scalable oversight}. Here, it is assumed that our true preferences are encoded by a function $f: \X \rightarrow \Y$ that is hard for us to compute, but easy for an AI. For example, suppose $x \in \X$ encodes the current state of the world and $y \in \Y$ is the desired action. If this involves long-term planning then $f$ can be hard to compute for us. There are many reasons for why an AI can be more powerful than us. The scalable oversight literature assumes that the reason is largely computational. For example, early works assumed that $f$ was a PSPACE-complete problem and that an AI was able to solve it~\citep{irving2018debate}, while a human was only able to solve problems in P. Some more recent works assume other kinds of computational superiority~\citep{browncohen2024debate,browncohen2025prover}. Since we are not able to compute $f$ ourselves, it is unclear how we can: a) design a training process that teaches a model to compute $f$, and b) check that a trained AI is indeed computing $f$. Scalable oversight aims to provide an answer. However, the currently existing literature has mostly focused on (b) while largely ignoring (a). 

That (b) might be achievable can be seen using the following analogy that appeared in~\cite{irving2018debate}. Suppose $x$ represents a SAT instance and $y$ denotes whether it is satisfiable. To prove to the human that $y=\mathrm{true}$, the AI just needs to provide the satisfying assignment along with $y$. In general, the scalable oversight literature assumes that $f$ sits within some level of the polynomial hierarchy, thus providing a means to certify whether AI's solution is correct. For example, suppose $f$ lies in the intersection of the two classes at level $2n$ of the polynomial hierarchy (at level 1, these are NP and co-NP).\footnote{Formally, this is the class $\Sigma_{2n}^{\mathrm{P}} \cap \Pi_{2n}^{\mathrm{P}}$.} Then there exists a polynomial time verifier $v$ such that for all $x$, the correct answer is $y$ iff $\exists c_1 \forall c_2 \ldots \exists c_{2n-1} \forall c_{2n}\, v(x,y,c_1,\ldots ,c_{2n}) = \mathrm{true}$. This motivates the \emph{debate} protocol~\citep{irving2018debate,browncohen2024debate}, where the AI is thought of as a pair of \emph{debaters} (also implemented as AI models). Given $x$, the \dt proposes a $y$ and tries to defend it. The \df argues for the opposite label. The two debaters  generate alternating arguments $c_1,c_2,\ldots , c_{2n-1},c_{2n}$, and the human judge $v$ judges whether the debate transcript proves that the answer to $x$ is $y$. The \dt's goal is to make the verdict `yes' and the \df's goal is the opposite. Due to the definition of this complexity class, one has the guarantee that the \dt has a winning strategy for a given $x$ if and only if its proposed label is equal to $f(x)$. If we allow the number of rounds of the debate to grow polynomially with the size of $x$, the problems whose answers can be certified this way are exactly those in PSPACE, and thus it is tempting to say that a human judge could potentially supervise models that solve problems in PSPACE (which was indeed the motivation presented in the first paper on debate in \citealp{irving2018debate}).


However, as mentioned above, the current literature ignores the important part (a), namely, how do we ensure that a model \emph{learns} the correct function $f$. Just because the existence of \dt's winning strategy is equivalent to the \dt returning the correct label, it does not imply that the winning strategy can be found. Moreover, it is likely that the advantage that an AI might have over humans is not as much about computation as about the ability to learn facts about the universe that humans haven't learned. Indeed, according to the extended Church-Turing thesis~\citep{vergis1986complexity, aaronson2005physical}, every phyically realizable computer can be \emph{efficiently} simulated by a (quantum) Turing machine. Thus, a computational intelligence, no matter how ``superhuman'' it gets, is unlikely to be solving worst case instances of NP-hard problems, let alone PSPACE-hard problems. If an AI is seemingly able to solve NP-hard problems (such as AlphaFold~\citep{jumper2021alphafold}), it is likely because it has \emph{learned} something about the distribution from which the universe is creating the input instances. Thus, we claim that learning should be at the centre of any investigation into scalable oversight. 

% And as the vast literature on learning theory has shown, learning is not always easy, even when humans \emph{are} able to provide the correct labels $y$. Indeed, despite the successes of neural networks, gradient descent fails to train them on the class of parity functions~\citep{shalevshwartz2017failures} \ruth{I remember that paper, but somehow i have a vague memory that parity is not an issue anymore.. or was it a formal hardness result? i guess i don't remember the paper really..}~\vinayak{Yeah this is not the best reference for this, but there are more recent empirical papers on this and my understanding is that it fails still for large enough dimensions. I need to cite those recent papers.} as predicted by statistical query lower bounds~\citep{kearns1998statistical,blum1994statistical}. Thus learning should be at the centre of any investigation into scalable oversight.

\subsection{Our results}
We model debate as an online learning problem. Nature produces a sequence of $x_t$'s, and at each round $t$, two debaters produce a sequence $\sc_t$ of certificates. The human is modeled by a verifier $v^*$ as above. Thus, given $x_t$, a proposed $\hy_t$, and a debate transcript $\sc_t$, $v^*$ can decide whether the transcript proves that $\hy_t$ is the correct label for $x_t$ and provides this feedback to the learner. Each verifier $v$ induces a ground truth hypothesis $h$. The debaters need to learn the true $h^*$. They don't know the true $v^*$ (because otherwise they would also know $h^*$) but know a class $\V$ that $v^*$ lies in. A \emph{mistake} is when the label $\hy_t$ proposed by the \dt is different from the true label $h^*(x_t)$. In the \emph{realizable} case, we assume that  the feedback is consistent with the true $v^*$ at each round.

Modeled this way, the learning problem becomes explicit and already opens up intriguing questions: First, what objective should adversarial debaters optimize? We show that the loss objective from previous work is not as useful in a learning setting, and propose a regret objective instead. Second, is it always better for the debaters to play adversarially? For example, can there be a case where the learning problem is solved more quickly if the two debaters cooperate? We answer this in the affirmative, where we show that there exist instances of the debate problem where learning is impossible when the debaters act adversarially to each other, but can be done with a constant number of mistakes if they cooperate. These results are presented in Section~\ref{sec:oppositional}.

Motivated by this example, we define the \emph{cooperative} version of the problem,
% \ruth{the previous paragraph, last sentence, already mentioned the cooperative case, so we need to restructure..}~\vinayak{the previous paragraph was using the English word cooperate to motivate that we need to define a formal framework to capture this phenomenon. This paragraph is talking about that formal framework.}\ruth{ok :)}, 
in which the debaters work towards the joint goal of using access to $v$ in the most efficient way in order to learn the ground truth. We present a characterization of the optimal mistake bound by (deterministic) cooperative debaters in the realizable setting (Section~\ref{sec:cooperative-case}), using a new combinatorial dimension that we define, termed \emph{\debatedimname}. A significant challenge in this setting is that in contrast with standard online learning protocols, the learner never observes the true labels that it is trying to predict. 

Another criticism of the debate protocol is that it assumes the infallibility of humans~\citep{irving2019social,browncohen2024debate,parrish2022single,parrish2022two}. It is unclear that a human can implement such a $v$ that aligns with the ground truth on all $x$. Thus, we also consider an unrealizable version of the problem, in which we assume that $v$'s output can be arbitrarily corrupted during some of the rounds. We call these rounds \emph{misspecified}, and show that for every constant $d > 1$, there are instances of the problem such that in the realizable setting, the optimal deterministic mistake bound is $d$, while if there are $k$ misspecified rounds, 
%%%%% start text for randomized lower bound 
then every learner makes at least $\Omega(dk)$ mistakes. This contrasts with Online Learning, in which any hypothesis class with a finite Littlestone dimension $L$ admits a ``no-regret'' algorithm for the $k$-non-realizable setting with a mistake upper bound of $k +O(\sqrt{k\cdot L} + L)$, implying that the prediction error rate converges to the misspecification rate \citep{FilmusHaMeMo23,BenDavidPaSh09,AlonBeDaMoNaYo21}. 
%%%%%% end text for randomized lower bound, deterministic is below
\iffalse 
then every deterministic learner makes at least $\Omega(dk)$ mistakes. This contrasts with Online Learning, in which any hypothesis class with a finite Littlestone dimension $L$ admits a randomized ``no-regret'' algorithm for the $k$-non-realizable setting with a mistake upper bound of $k +O(\sqrt{k\cdot L} + L)$, and a deterministic algorithm with at most twice as many mistakes, implying that the prediction error rate converges to at most twice the misspecification rate \citep{FilmusHaMeMo23,BenDavidPaSh09,AlonBeDaMoNaYo21}. 
\fi
Our lower bound for debate learning shows that debate learning is significantly more  sensitive to misspecification errors: even a small rate of misspecification errors can induce an unreasonably large prediction error, and the error rate may never approach the misspecification rate. 
%\sivan{Tosca, if you manage to prove a lower bound for a randomized adversary, the text above should be updated including the compoarison to online learning}

For some of our results, the proofs are deferred to the appendix.

\subsection{Related work}
There is very little previous work that analyzes learning in the
setting of debate. Perhaps the two that are closest to our work
are~\citet{chen2023playing} and~\citet{buzaglo2025hidden}. Both model
debate as a repeated game between the two debaters with a fixed payoff
matrix. Since the space of arguments is too large to enumerate, the
debaters access this matrix through an optimization oracle, and both
papers give efficient no-regret strategies under this access. By
classic results, the time-averaged play then converges to an
equilibrium~\citep{freund1999adaptive,hart2000simple}. In this model,
the debaters never need to learn the verifier, and it is assumed that the equilibrium answer is correct.
In our model, the verifier is unknown, each round reveals a single
verdict on a new input, and we measure success by the number of
mistakes with respect to the ground truth.



% \section{Related Work}

% There is very little previous work that analyzes learning in the
setting of debate. Perhaps the two that are closest to our work
are~\citet{chen2023playing} and~\citet{buzaglo2025hidden}. Both model
debate as a repeated game between the two debaters with a fixed payoff
matrix. Since the space of arguments is too large to enumerate, the
debaters access this matrix through an optimization oracle, and both
papers give efficient no-regret strategies under this access. By
classic results, the time-averaged play then converges to an
equilibrium~\citep{freund1999adaptive,hart2000simple}. In this model,
the debaters never need to learn the verifier, and it is assumed that the equilibrium answer is correct.
In our model, the verifier is unknown, each round reveals a single
verdict on a new input, and we measure success by the number of
mistakes with respect to the ground truth.


\section{Formal Setup}\label{sec:formal}
% We let $\X$ denote a \emph{domain}, $\Y = \{0,1\}$ a binary \emph{label space}, and $\C = \{0,1\}$ a binary \emph{certificate space}.  
We let $\X$ denote a \emph{domain}, $\Y = \{0,1\}$ a binary \emph{label space}, and a finite $\C$ a \emph{certificate space}.  
A \emph{hypothesis} is a function $h:\X\to\Y$, and we use $\H$ to denote a \emph{hypothesis class}.
A \emph{debate of length} $n$ is a sequence $(\c_1, \c_2, \ldots, \c_{2n})$ of $2n$ certificates $c_i\in\C$. We think of these certificates as being produced alternately by two debaters. That is, we think of the certificates $c_i$ of odd index as being produced by the first debater, which we will call the \dt, and the certificates $c_i$ of even index by the second debater, which we will refer to as the \df. To clarify the role of the arguments in the debate, we usually will use the notation $(\ct^1, \cf^1, \ct^2, \cf^2, \ldots,  \ct^n, \cf^n)$ for a debate of length $n$

The debate is moderated through a verifier, which takes in a domain point and a debate of length $n$ and produces a label for $x$ given the arguments in the debate. 
The following definition specifies the relationship between hypotheses and verifiers in our framework.

\begin{definition}
Let $h:\X\to\Y$ be a hypothesis. We say that verifier $v: \X\times \C^{2n} \to \Y$ \emph{validates} hypothesis $h$ through debate of length $n$ if for all $x\in\X$ there exists a unique $y\in \Y$ such that
\begin{align*}
\exists\ct^1\in\C\,\forall\cf^1\in\C\,\exists\ct^2\in\C\,\forall \cf^2\in\C\ldots\,\exists\ct^n\in\C\,\forall\cf^n\in\C,
v(x, \ct^1, \cf^1, \ct^2, \cf^2, \ldots,  \ct^n, \cf^n) = y,
\end{align*}
and $h(x) = y$ for this $y$. 
% \begin{align*}
% ~\exists~\ct^1\in\C~\forall~\cf^1\in\C~\exists~\ct^2\in\C~\forall~\cf^2\in\C\ldots~\exists~\ct^n\in\C ~\forall~\cf^n\in\C ~:~
%  v(x, \ct^1, \cf^1, \ct^2, \cf^2, \ldots,  \ct^n, \cf^n) = y
%  % \\ \iff h(x) = y
% \end{align*}
%is equivalent to $h(x) = y$.
% \vinayak{
% This ensures that the situation where there is a valid debate transcript for both labels does not happen.
% }
In this case, we also say that $v$ is a \emph{verifier} for hypothesis $h$. We call a verifier \emph{valid} if there exists a hypothesis $h$ that is validated by $v$.
\end{definition}

Note that a given valid verifier $v$ validates exactly one hypothesis $h$ in the above sense. We will use notation $h_v$ for this hypothesis when we want to make this relationship between verifier and its hypothesis explicit. On the other hand, a given hypothesis $h$ may be validated by multiple verifiers $v$.
% \ruth{This is still true? Does it matter? Or should we assume $\V$ is pruned in a way that there is a one to one correspondence between $\H$ and $\V$?}\vinayak{Yes this is true, and I think this is kind of important, so we should keep it this way.}\sivan{I agree with that this is important to keep}
% ~\vinayak{Oh I see. I have been thinking of verifiers as uniquely defining a hypothesis. Given an $x$ and a $v$, I think only one unique label should meet the condition above, and that label will define a unique hypothesis. I think Tosca is saying the same thing below.}. 
We call a given set of valid verifiers a \emph{verifiers class}, and denote it $\V$. Further, we let $\H_\V$ denote the hypothesis class of all functions that are validated by verifiers $v\in\V$. If clear from context, we may simply write $\H$ for $\H_\V$.
For a hypothesis class $\H$ we say that $\H$ is \emph{supported by verifier class} $\V$ if $\H = \H_\V$.

% \tosca{I think one thing that is still missing here is that we want the label that is certified to be unique. In the current definition there exists a hypothesis that verfies \emph{all hypotheses}. I suggest we add "soundness" of a verifier somewhere (i.e. it for every $x$ there exists no more than one label $y$ that $v$ verfifies)}

\paragraph{Learning via debate} The goal of \emph{learning via debate} is to recover an underlying \emph{true hypothesis}, which we denote by $h^*$. We assume that the two debaters interact with a verifier $v^*\in\V$ that validates the true hypothesis $h^*$ (and thus $h^*\in\H = \H_\V$).
In the \emph{realizable case} all label feedback is provided by the verifier $v^*$.
In the \emph{$k$-non-realizable} case, the label feedback in at most $k$ rounds may be different from that of $v^*$. We call these rounds \emph{misspecified}. Of course, debaters \dt and \df do not know the identity of $h^*$ nor the identity of $v^*\in\V$. They only have access to the verifier class $\V$, and thereby also to its induced hypothesis class $\H_\V$.

We first give the debate protocol for the realizable setting. 
\paragraph{The debate protocol:}
\begin{itemize}
\item For rounds $t = 1\ldots T$ repeat:
\begin{enumerate}
 \item The environment poses an $x_t\in\X$
 \item Debater \dt hypothesizes label $\hy_t$ (also called ``prediction'' or ``predicted label'')
 \item For $i = 1 \ldots n$ repeat:
 \begin{itemize}
  \item Debater \dt brings argument $\ct^i_t$ (in support of label $\hy_t$)
  \item Debater \df brings argument $\cf^i_t$ (aiming to falsify label $\hy_t$)
 \end{itemize}
Let $\sc_t=(\ct^1_t,\cf^1_t,\ldots,\ct^n_t,\cf^n_t)$.
 \item The debaters receive $\vy_t=v^*(x_t,\sc_t)$. 
 \item %~\vinayak{The environment adversarially decides whether round $t$ is misspecified.  The debaters receive $\vy_t=v^*(x_t,\sc_t)$ if the round is not misspecified and $\vy_t=1-v^*(x_t,\sc_t)$ if it is. They are not told whether the round is misspecified.}\\
       The environment records (but does not reveal) $\ty_t = h^*(x_t)$ 
\end{enumerate}
\end{itemize}

A run of the debate protocol is evaluated through the \emph{number of mistakes} in predicting the label $\ty_t$. For a fixed interaction sequence, denote the number of mistakes by
%\begin{align*}
$\Mi(T) = \sum_{t = 1}^T \indct[\hy_t \neq \ty_t].$
%\end{align*}





A standard approach for using debate as a learning strategy considers the two debaters as playing the zero-sum game of winning the argument~\citep{irving2018debate,chen2023playing,browncohen2024debate,browncohen2025prover}.
Standard formalizations \citep{irving2018debate,browncohen2024debate} define a loss for each debater.\footnote{Some works let the verifier output a score rather than a binary verdict, measuring how strongly the debate transcript supports the proposed label~\citep{chen2023playing,arnesen2024training,browncohen2025prover}; we use binary verdicts throughout.} Debater \dt aims to defend its hypothesized label $\hy_t$ and is therefore evaluated through the following loss: $\tLo(T) = \sum_{t = 1}^T \indct[\vy_t \neq \hy_t].$ 
Debater \df, on the other hand, is evaluated against its aim to falsify the hypothesis that the first debater proposed, by setting its certificates in such a way that the verifier will conclude with a label $y_t$ that is not the hypothesized label $\hy_t$. 
Therefore, \df is evaluated through the opposite loss: 
%\begin{align*}
 $\fLo(T) = \sum_{t = 1}^T \indct[\vy_t = \hy_t] = T - \tLo(T).$
%\end{align*}
If the learning problem is ignored, so that $v^*$ is known to both debaters, then a \dt optimizing $\tLo(T)$ with respect to a \df that optimizes $\fLo(T)$ leads to \dt always proposing the true label of $h^*$ verified by $v^*$, with certificates that cannot be refuted. Moreover, this leads to $\Mi(T) = \tLo(T)$, motivating this optimization objective. We discuss below the setting of two debaters that optimize their objectives separately, which we term the \emph{oppositional setting}. 

We can also consider \dt and \df that can be cooperative, with the joint goal to minimize $\Mi(T)$. 
In this case, \df provides its certificates in good faith ``for the sake of the argument'' so as to lead the verifier to the correct conclusions, having taken into account relevant potential counter-arguments. In this case, the two debaters can optimize their behavior together and can be thought of as a single entity that we also sometimes refer to as \emph{the learner} in this case. We study this learning setting below.


%In principle, this setting, in which each debater attempts to optimize its own loss, can lead \dt to propose hypothesized labels $\hy_t$ that it knows it can argue for, even though it is already aware that this does not correspond to the underlying true label $y_t$. 
%\vinayak{Is this assumption also standard in the debate literature? if so we need references (sivan)}~\vinayak{I am not sure if there is an assumption here. I think Ruth is just trying to say that minimizing the debater's loss isn't the same as minimizing mistakes.} 

%We assume that each debater is unaware of the behavior of the other debater, and therefore runs its strategy assuming that the other debater can be adversarial to it, providing worst-case sequences with respect to its loss objective. \sivan{depends on lemma: This is optimal, in the sense if \dt and \df were each optimal against their opponent, then their worst-case loss would be equal to the worst-case loss that can be obtained by a debater that assumes that the other debater is adversarial (see Lemma~\ref{lem:equilibrium} in the appendix).\footnote{Note, however, that such an equilibrium does not necessarily exist for deterministic debaters \citep{vonneumann1928theorie}.}We call this the \emph{oppositional setting}. }




% partition $C= C_0 \cup C_1$
% if $h^* = h_0$, then winning certificate $\ct\in C_0 for \dt$
% if $h^* = h_1$, then winning certificate $\ct\in C_1 for \dt$

% So winning strategy for \dt will always choose $\ct$ from $C_{\hy}$.

% There is $\cf \in C$ such that for every $\ct \in C_1$ and every $x\in \X$ and every $v\in \V$
% $v(x,\ct,\cf) = 1$

% There is $\cf \in C$ such that for every $\ct \in C_0$ and every $x\in \X$ and every $v\in \V$
% $v(x,\ct,\cf) = 0$











\section{Cooperative Debaters}
\label{sec:cooperative-case}
In this section, we start by exploring learning via debate using cooperative debaters. As mentioned above, we will call them jointly \emph{the learner} in this case. To analyze which classes $\V$ of verifiers allow for learning with finite mistake bound, we will define a tree structure that represents potential runs of the debate protocol for debate length $n$.

In our tree construction, edges correspond to all possible learner actions 
and nodes correspond to worst-case actions of the environment. 
This type of construction has been previously used in \cite{oz2026discriminative} for a different interactive learning protocol, Discriminative Feature Feedback. Nonetheless, as we will see below, our definition of dimension based on the tree will be quite different from standard ones, due to the special challenges of the debate protocol.

In our case, the learner's action in each round includes the label $\hy$ that the \dt poses, as well as the debate $\sc$ of length $n$ that then ensues between the two debaters.
The environment's actions include a verifier evaluation $\vy$ for some verifier consistent with the previous rounds, as well as the  domain instance that the environment poses for the next round.
The root node only records the instance that the protocol is initiated with.



\begin{definition}[Debate Tree]
Given a debate length $n$ and a verifier class $\V$, we define a \emph{length-$n$-debate tree for class $\V$} to be a directed tree $\T$ with the following properties:
\begin{enumerate}
    \item Every node $\alpha$ in $\T$ is a pair of a domain point and a label $\alpha=(x_\alpha, \vy_\alpha)$.
    \item The root node is denoted by $\rho$. Its domain point is denoted by $x_\rho$, and its label is $\vy_\rho = \sqcup$, where $\sqcup\notin\Y$. \item The root node $\rho$ is the only node in $\T$ with associated label $\vy_\rho \notin\Y$: for all nodes $\alpha\neq \rho$, we have $\vy_\alpha\in \Y$.
    \item Edges $(\alpha, \beta)$ in $\T$ are directed away from the root node, and annotated with a label $\ell^\alpha_\beta = (\hy^\alpha_\beta, \sc^\alpha_\beta)$, where $\sc^\alpha_\beta = (\ct^1_{\alpha\beta}, \cf^1_{\alpha\beta}, \ct^2_{\alpha\beta}, \cf^2_{\alpha\beta}, \ldots,  \ct^n_{\alpha\beta}, \cf^n_{\alpha\beta})$ is a debate of length $n$.
    % We will denote edges by triples
    % \[
    % (\alpha, \ell, \beta) = ((x_\alpha, y_\alpha), (\hy^\alpha_\beta, \sc^\alpha_\beta),  (x_\beta, y_\beta)).
    % \]
    \item For each non-leaf node $\alpha$ and each possible edge annotation $(\hy, \sc)$, there is exactly one child node $\beta$ to $\alpha$ with $\ell^\alpha_\beta = (\hy, \sc)$
    \item All directed paths in $\T$ are finite, and consistent with a verifier $v\in\V$ (in the sense defined below).
\end{enumerate}
We let $\Delta_n(\V)$ denote the set of all length-$n$-debate trees $\T$ for verifier class $\V$.
\end{definition}


Let $(\alpha,\beta)$ be an edge in the tree $\T$ and let $v\in\V$ be a verifier. We say \emph{$v$ is consistent with edge $(\alpha, \beta)$} if
\[
v(x_\alpha, \sc^\alpha_\beta) = v(x_\alpha, \ct^1_{\alpha\beta}, \cf^1_{\alpha\beta}, \ct^2_{\alpha\beta}, \cf^2_{\alpha\beta}, \ldots,  \ct^n_{\alpha\beta}, \cf^n_{\alpha\beta}) = \vy_\beta.
\]
Let $p$ be a directed path (down the tree) in $\T$. We say that the path $p$ is consistent with verifier $v$ if all edges in $p$ are consistent with $v$.

To analyze learnability, we will define \emph{version spaces} of verifiers that agree with paths in the debate tree.
 We define the \emph{verifier version space of path $p$} to be the set
\[
 \V_p = \{ v\in\V ~\mid~ v \text{ is consistent with all edges in }p\}.
\]
Recall that each verifier $v$ induces one hypothesis $h$. Therefore, a verifier version space $\V_p$ induces a corresponding hypothesis version space $\H_p$. Note also that more than one verifier may induce the same hypothesis.

For a descendant node $\beta$ of a node $\alpha$ in $\T$, we let $p^\alpha_\beta$ denote the directed path from $\alpha$ to $\beta$.  In case $\alpha =\rho$ is the root node, we may simply use the notation $p_\beta = p^\rho_\beta$ to denote the path from root to node $\beta$. We will let $P^\alpha$ denote the set of all root-to-leaf paths in $\T$ that start at $\alpha$.

For version spaces, we may use the shorthand notations $\V^\alpha_\beta = \V_{p^\alpha_\beta}$ and $\V_\beta = \V_{p_\beta}$. That is $\V_\beta$ is the verifier version space at a node $\beta$ in the tree: $\V_\beta$ contains exactly those verifiers that are consistent with the debate protocol as run from the beginning $\rho$ to node $\beta$ in the tree. Analogously, we use $\H^\alpha_\beta$ and $\H_\beta$ for the corresponding induced hypothesis version spaces.

For any hypothesis $h:\X\to\Y$, we let $\mi(h, p )$ denote the number of disagreements between $h$ and the hypothesized labels on the edges of path $p$. That is $\mi(h, p) = \sum_{(\alpha, \beta)\in p}\indct[h(x_\alpha)\neq \hy^\alpha_\beta]$.
Note that if a verifier $v$ is consistent with path $p$ and $h_v$ is the hypothesis corresponding to $v$, then  $\mi(h_v, p)$ corresponds to the number of mistakes the debaters incurred assuming that candidate verifier $v$ is the true $v^*$.

% \sivan{I don't think that's true - the verifier can return labels that are not true if it gets bad arguments}\ruth{yes, that's why we're comparing $h$ to the the labels $\hy$, and don't make any conclusions from the potentially bad labels $\vy$, other than saying $h_v$ is ``still in the game'', and $\mi(h_v, p)$ would be the number of incurred mistakes if $v$ was the true verifier; so statement seems fine to me}\sivan{I understood "mistakes the debaters incurred with respect to candidate verifier $v$" as rounds in which $\hy \neq \vy$. I think you meant something else. I reworded to clear this up.}\ruth{sounds good}

For a tree $\T$ and node $\alpha$ in $\T$, we let $\T^\alpha$ denote the subtree rooted at $\alpha$ in $\T$. Recall that $P^\alpha$ is the set of all root-to-leaf paths in $\T$ that are starting at $\alpha$.
We define the \emph{depth} of subtree $\T^\alpha$ with respect to verifier class $\V$ as follows:
\[
\depth(\T^\alpha, \V) = \min_{q\in P^\alpha}~ \max_{h\in \H_\alpha \cap \H_q}~  ( \mi(h, p_\alpha) + \mi(h, q) ).
\]
% \[
% \depth(\T^\alpha, \V) = \min_{p\in P^\alpha}~ \max_{\beta\in p}~ \max_{h\in \H_\alpha \cap \H^\alpha_\beta}~  ( \mi(h, p_\alpha) + \mi(h, p) )
% \]

Note that the depth of a subtree rooted at $\alpha$ counts the number of (past and future) mistakes that the debate could incur with verifiers that are consistent with the path up to $\alpha$.
The use of past mistakes in the depth of a subtree is crucial: because the learner (the debaters) never directly finds out the true  labels $\ty_t$, it must keep track of all consistent hypotheses and its mistakes with respect to them throughout the run.
The depth of the whole tree $\T$ is simply its depth at root node $\rho$, that is $\depth(\T) = \depth(\T^\rho)$. Note that at the root node, there is no past, and thus only future mistakes are counted:
\[
\depth(\T,\V) = \depth(\T^\rho,\V) = \min_{q\in P^\rho}~ \max_{h\in \H_\V \cap \H_q}~  ( \mi(h, p_\rho) + \mi(h, q) )  = \min_{q\in P^\rho}~ \max_{h\in \H_{q}}~    \mi(h, q).
\]

% \[
% \depth(\T,\V) =  \min_{p\in P^\rho}~ \max_{\beta\in p}~ \max_{h\in \H_\V \cap \H_\beta}~  ( \mi(h, p_\alpha) + \mi(h, p) )  = \min_{p\in P^\rho}~ \max_{\beta\in p}~ \max_{h\in \H_\beta}~    \mi(h, p)
% \]
%
% \ruth{there is an issue right now, since we cannot allow the minimum just choose a short path... need to take min over all ``maximally consistent'' paths.. so will need more notation..}
With this, we are ready to define the \emph{\debatedimname} for debate length $n$ of a verifier class:
\begin{definition}[\debatedimname]
 Let $\V$ be a verifier class for domain $\X$. Let $\Y$ be a label space, $\C$ be a certificate space, and $n\in\naturals$ be a debate length. The \emph{length-$n$-\debatedimname} $\dldim(\V, n)$ of $\V$ is the largest possible depth of a debate tree for debate length $n$:
 \[
  \dldim(\V, n) = \sup_{\T\in\Delta_n(\V)} \depth(\T,\V).
 \]
If the debate length $n$ is fixed (or clear from context) we will also use the shorthand notation $\dldim(\V) =   \dldim(\V, n)$.
\end{definition}
We will show that the \debatedimname characterizes when a hypothesis class $\H$ that is supported by a verifier class $\V$ is learnable via deterministic debate with a finite mistake bound. 

\begin{theorem}
Let $\H\subseteq \Y^\X$ be some hypothesis class. Then $\H$ is learnable from length-$n$-debate with finite mistake bound by a deterministic learner if and only if the learner has access to a verifier class $\V$ with $\H = \H_\V$ and $\dldim(\V, n) <\infty$.
Moreover, the optimal mistake bound for learning $\H$ through access to $\V$ is exactly $\dldim(\V, n)$, the length-$n$-\debatedimname of $\V$.
\end{theorem}
% \ruth{Maybe, for a given hypothesis class, the dldim for all supporting verifier classes is equal? Not sure right now, so phrasing it as above..If we have the same dimension for all supporting verifier classes, we could also define the dldim(H); otherwise, we should define it as inf over all supporting verifier classes..need to think through}
% \sivan{I don't think this holds. You could have a verifier class that reveals H exactly using "hidden" values, while another class does not. I'm not sure if we need a dimention for H though}\ruth{ok, removing the comment}

The theorem follows directly from the subsequent two lemmas in this section. Lemma \ref{lemma:cooperative_deterministic_lowerbound} shows that our dimension lower bounds the optimal mistake bound, while Lemma \ref{lemma:cooperative_deterministic_upperbound} provides the corresponding upper bound.
We start by showing that the \debatedimname induces a lower bound on the number of mistakes that the environment can enforce.

\begin{lemma}\label{lemma:cooperative_deterministic_lowerbound}
% Let $\C$ be a finite certificate class. 
Let $\H\subseteq \Y^\X$ be some hypothesis class, for which there exists a verifier class $\V$ with $\H = \H_\V$ and $\dldim(\V, n) = d \in \naturals$. Then, for any deterministic choices of the debaters in the debate protocol with verifier class $\V$, the environment can enforce $d$ mistakes within finitely many rounds.
% \begin{align*}
%  \Mi(T) = \sum_{t = 1}^T \indct[\hy_t \neq \ty_t] \geq d
% \end{align*}
\end{lemma}

\begin{proof}
Assume $\dldim(\V, n) = d$. 
Since $d\in\naturals$ is finite, there exists a debate tree $\T$ with $\depth(\T,\V) = d$. 
From the definition of depth
\[
\depth(\T,\V) =  \min_{p\in P^\rho}~  \max_{h\in \H_p}~ \mi(h, p)
\]
it follows that for any $p \in P^\rho$, there is some $h$ that induces $d$ mistakes. Let $t_p$ denote the round after which this $h$ has induced $d$ mistakes in path $p$. Note that every $p$ in $\T$ is finite, and the fan-out of the tree $\T$ is finite (since $\C$ is finite). Hence, 
$\max_{p\in P^\rho} t_p$ is some $b \in \naturals$. 
It follows that for any sequence of outputs chosen by the learner, the environment has verifier responses and a consistent choice of $h$ that enforce $d$ mistakes in at most $b$ steps. 
\end{proof}
\iffalse
implies that there is uniform upper bound $b < \infty$, such that the $d$ mistakes that are possible on every path $p$ in the tree (ie for any choice of the debaters), can be enforced by the environment within $b$ steps. \ruth{argument could be expanded, assume there doesn't exist that uniform upper bound $b$, then at every step there exists a choice where the mistakes can be postponed arbitrarily long, and that leads to an infinite path in the tree}
\fi
% \ruth{actually might not be sufficient when the certificate space in infinite? yes.. we need bounded fan-out for this argument..}, the definition of depth
% \[
% \depth(\T,\V) =  \min_{p\in P^\rho}~ \max_{\beta\in p}~ \max_{h\in \H_\beta}~ \mi(h, p)
% \]
% implies that 
% \ruth{For trees: there does not exist a tree with finite fan-out and arbitrarily long paths, that does not also have an infinite path. But with infinite fan-out this is not true, one can have a node with finite paths of every length starting at that node... }
%\end{proof}

For the upper bound, we define the Standard Optimal Algorithm (SOA) for the length-$n$-debate protocol. Assume that the debate protocol is now in round $t$, and let $p_t$ be the path (of a debate tree) recording the previous rounds. Letting $x_t,\hy_t,\sc_t,\vy_t$ be the record of round $t$, we define $p_t$ to include $x_t$ but not $\hy_t,\sc_t,\vy_t$.
% \ruth{should this be $\vy_{t+1}$? or do we let the indices of the $x$ and $\vy$ in a node in the tree be different?}\sivan{yes they are different, $\vy_t$ corresponds to the response to $\hy_t$ and the same node includes the next example which is $x_{t+1}$. }


Then the current verifier version space is $\V_{p_t}$. 
An optimal learner needs to take the action $(\hy, \sc)$ that minimizes the maximal number of mistakes the environment will be able to enforce. The environment can enforce mistakes by combining past mistakes that happened on $p_t$ with future mistakes that may happen in a debate tree from this point on \emph{through the same function $h$}.
Thus, the learner chooses a pair $(\hy, \sc)$ in $\argmin_{\hy, \sc}Q_t(\hy, \sc)$ for

% Thus, our learner chooses $(\hy, \sc)$ that minimizes
% \sout{
% %\[
% $Q = \max_{\vy\in\{0,1\}} ~\sup_{\T\in\Delta_n(\V_{(q,\hy, \sc, \vy)})} ~\min_{p\in\T^\rho} ~\max_{h\in \H_p\cap \H_{(q,\hy, \sc, \vy)}} (\mi(h, p) + \mi(h,q)),$
% %\]
% }
% \sivan{
\[
Q_t(\hy, \sc) = \max_{\vy\in\{0,1\}} ~\sup_{\T\in\Delta_n(\V_{(p_t,\hy, \sc, \vy)})} ~\min_{q\in\T^\rho} ~\max_{h\in \H_{(p_t,\hy, \sc, \vy)}\cap \H_q} (\mi(h, p_t) + \mi(h,q) +\indct[h(x_t) \neq \hy]),
\]
% }
where we use the notation $(p_t,\hy, \sc, \vy)$ for the path $p_t$ expanded with learner actions $(\hy, \sc)$ and environment response $\vy$. 
Note that for the version spaces defined through this expanded path, the domain point component $x_{t+1}$ of the  environment response is not relevant, and we thus omit it here to not overload the notation.

\begin{lemma}\label{lemma:cooperative_deterministic_upperbound}
 Let $\H\subseteq \Y^\X$ be some hypothesis class, for which there exists a verifier class $\V$ with $\H = \H_\V$ and $\dldim(\V, n) \in \naturals$. Then, at every step $t$ of the SOA learner above, $\min_{\hy,\sc} Q_t(\hy,\sc) \leq \dldim(\V, n)$. 
  Moreover, $\dldim(\V, n)$ is also an upper bound on the total number of mistakes that the SOA learner makes during the run.
\end{lemma}


% \ruth{The following still needs to be shown, but should be ok:}\sivan{are we using this lemma anywhere? we can remove if we don't}
% \begin{lemma}
% Let $\T$ be a debate tree and let $\alpha$ be a node in $\T$. Then
% \[
% \depth(\T^\alpha, \V) \leq \depth(\T, \V)
% \]
% \end{lemma}



\section{Oppositional Debaters}
\label{sec:oppositional}

We now turn to consider the oppositional setting. As described in Section~\ref{sec:formal}, the standard approach (that does not involve learning) considers debaters that optimize the worst case $\tLo$ and $\fLo$ with respect to a worst-case opposing debater. However, we show here that in a learning setting, an approach minimizing a notion of \emph{regret} is more suitable. 

\subsection{Using regret to upper bound prediction mistakes}

Given $v^*$, there is a strategy for \dt that obtains $\tLo(T) = 0$, since for any $x_t$, it can choose $\ct^1,\ldots,\ct^n$ that satisfy the validation requirement of $v^*$ for any $\cf^1,\ldots,\cf^n$. We call a round in which \dt selected $\hy_t  = h_{v^*}(x_t)$ and such certificates an \emph{irrefutable round}. For \df, even given $v^*$, it cannot avoid incurring a loss in irrefutable rounds. Therefore, for \df a notion of regret instead of loss is relevant even in a realizable setting: Let $U(T)$ be the number of irrefutable rounds in a given interaction. The smallest possible loss for \df is therefore $U(T)$. Thus, its regret is defined as $\fRe(T) = \fLo(T) - U(T)$. For \dt in the realizable setting, its regret is equal to its loss: 
$\tRe(T) = \tLo(T)$. 


The following claim shows that the total regret of the debaters upper bounds the total mistakes of the debate protocol. 


% The value of the overall debate that the two debaters produce is measured by assessing whether the sequence of produced arguments/certificates actually lead the verifier to conclude the correct label $\ty_t$ in each round.

% \sivan{testing:

% Let $L$ be the maximum total loss of \dt with respect to worst case environment and \df. Then in all runs, for the optimal \df, $\Mi(T) \leq L$. 

% *** can we show an example where $\Mi(T) > L$ for some debater 2.

% if there is a response of debater 2 that would reveal $v^*$, then if debater 2 wouldn't provide it, then debater 1 will not know v* and make more prediction mistakes in the end. 



% Proof: consider a case where \df knows $v^*$. So whenever \dt makes a refutable argument, \df refutes it and \dt suffers a loss of $1$. Against this debater, \dt has a loss of at most $L$. 

% -- doesn't work because a different \df may intentionally not refute to reduce \dt's information.


% Let Debater 2 have a learning strategy with worst-case regret gurarantee
% R_2(D_2) \leq d.

% then for every sequence (x_1, \hy_1, c_1, \tilde{y}_1), ... ,(x_n, \hy_n, c_n, \tilde{y}_n)

% Loss_{sequence}(D_2) \leq \min_{c_2^1...c_2^n\in C} Loss_{sequence}((c_2^1, ...c_2^n )) + d

% }


\begin{claim}\label{claim:regret}
In any run of the debate protocol,  $\Mi(T) \leq \tRe(T) + \fRe(T)$.
\end{claim}
\begin{proof}
In an irrefutable round,  \dt's choices force $v^*$ to output $\vy_t = \hy_t$, regardless of the choices of \df. Since $v^*$ validates $h^* = h_{v^*}$, this means that $\hy_t=h^*(x_t)$, so there are no prediction mistakes in these rounds. Hence, $\Mi(T) \leq T-U(T)$. In every round, exactly one of the debaters incurs a loss, so
$T =  \tLo(T)+\fLo(T)$. Therefore
\[
\Mi(T) \leq \tLo(T)+\fLo(T)-U(T) = \tRe(T) + \fRe(T).
\]
As claimed.
\end{proof}

This upper bound motivates new objectives for \dt and \df in the oppositional setting when learning is involved: if each debater minimizes their worst-case \emph{regret} with respect to a worst-case opposing debater, this upper bounds the worst-case number of debate protocol mistakes: if the worst-case regret for \dt over $T$ rounds is $R_P^T$ and the worst-case regret for \df  over $T$ rounds is $R_R^T$, then for any sequence, $\Mi(T) \leq R_P^T+ R_R^T$.



%This motivates the oppositional setting: If each debater selfishly succeeds to obtain a lower regret, it also reduces the prediction error of the protocol.\footnote{\sivan{(not sure yet about this comment, related to regret vs loss issue) Note in particular, that if \df wants to minimize its regret, it can do so by helping \dt identify the true $v^*$, since this would lead to $U(T) = T$ and $\fRe(T) = 0$. However, this may require \df to make assumptions on the behavior of \dt and not treat its input as fully adversarial.}}



Moreover, we show next that unlike the regret optimization approach, the loss optimization approach discussed above does not provide similar guarantees and can in fact completely fail. 

\subsection{Loss optimization vs.~regret optimization}

We show that even if each debater optimizes its worst-case loss against worst-case environment and opposing debater, this does not provide helpful guarantees for learning. 
First, we observe that for \dt, this optimization criterion is  meaningless, since all \dt strategies are optimal according to this criterion. 

\begin{claim}\label{claim:deb2}
    For any verifier class $\V$ and any number of rounds $T$, the worst-case optimal loss for $\df$ is $T$. 
\end{claim}
\begin{proof}
Consider a sequence $(x_t,\hy_t,\ct_t^1, \ldots, \ct_t^n)_{t \in [T]}$ that is consistent with some verifier $v^* \in \V$ such that for each $t$, $\forall \cf_t^1, \ldots, \cf_t^n, v^*(x_t, \ct_t^1,\cf_t^1, \ldots, \ct_t^n, \cf_t^n) = \hy_t$. Note that such $\hy_t$ and choices $\ct_t^1, \ldots, \ct_t^n$ exists for every $x \in \X$ according to the definition of a verifier. 
For this sequence, any response of \df has $\hy_t = \vy_t$. Hence, \df incurs a loss in every round, leading to $\fLo(T) = T$.
\end{proof}

In addition, in the next theorem we show that the loss of \dt cannot be used to upper bound the debate protocol mistakes $\Mi(T)$ for \emph{loss-optimizing} debaters, that is debaters that are worst-case optimal with respect to their loss as indicated above. In contrast, if the debaters are worst-case optimal with respect to their \emph{regret}, then according to the bound in Claim~\ref{claim:regret}, these regrets upper bound $\Mi(T)$. In the same example, we show that $\Mi(T)$ is large under the loss optimization scenario but small under the regret optimization scenario. 
\begin{theorem}
    There is a verifier class $\V$ such that loss-optimizing debaters lead to 
    $\ex[\Mi(T)] \geq T/4$ for all environment inputs, whereas $\ex[\tLo(T)] \leq 1.5$. On the other hand, regret-optimizing debaters lead to $\ex[\Mi(T)] \leq \ex[\tLo(T)] = 1.5$ for all environment inputs.
\end{theorem}
\begin{proof}
Let $\X = \{x\}$ and $\C = \{0,1\}$. Let $\H = \{h_0,h_1\}$, where $h_b(x) = b$ for  $b\in\{0,1\}$. Let $\V$ be the class of all verifiers $\{ v_{b,c} \mid b,c \in \{0,1\}\}$  such that $\forall \cf \in \C, v_{b,c}(x,c,\cf) = b$, and for $1-c$ we have $v_{b,c}(x,1-c,\cf) = \cf$ for $\cf \in \C$, as illustrated below.
\[
\begin{array}{c|cc}
 v_{b,c}(x,\ct,\cf) & \cf=0 & \cf=1\\
\hline
\ct=c & b & b\\
\ct = 1-c & 0 & 1
\end{array}
\]
Note that $v_{b,c}$ verifies $h_b$. 

First, consider an optimal \dt. Note that loss optimization and regret optimization are equivalent for \dt. Let $b^* \in \{0,1\},c^*\in \C$ such that $v^* = v_{b^*,c^*}$. Consider the following \dt strategy: It maintains a set of active pairs of label and certificate, initialized to $\{0,1\}\times \C$. In the first round, it selects $\hy,\ct$ uniformly at random from the set of active pairs. In each round, if it incurs a loss using some pair, it removes this pair from the active set and selects another pair for the next round uniformly at random. If it does not incur a loss using in a round, it continues to use the same pair in the next round. 

To see that this is worst-case optimal, observe that whenever \dt proposes some $\hy$ and $\ct$ for some $x \in \X$ and incurs a loss with some \df certificate $\cf$, this means that 
$\vy = v_{b^*,c^*}(x,\ct,\cf) \neq \hy$. This implies that $v_{b^*, c^*} \neq v_{\hy,\ct}$, since $v_{\hy,\ct}(x,\ct,\cf) = \hy$. Hence, this \dt strategy only incurs a loss at most once for every pair until selecting $b^*, c^*$. The expected number of pair selections before selecting this pair is $1.5$, hence the expected loss of \dt is at most $1.5$. This is worst-case optimal, since any strategy that does not replace a pair that has incurred a loss will incur the loss again in the next round against a worst-case \df, thus its worst-case loss cannot be smaller.

Moreover, let $j$ be the last round before $b^*,c^*$ is found, then using this \dt strategy we get $\Mi(T) \leq j$, since no prediction mistake is made after this round. However, the value of $j$ depends on the strategy of \df, as we now show.

First, consider the optimal \df according to loss minimization. As shown in Claim~\ref{claim:deb2}, all \df strategies are optimal with respect to worst-case loss minimization. In particular, consider the strategy that in every round, upon observing $(x_t, \ct_t, \hy_t)$, selects the response $\cf_t = \hy_t$. Under this strategy, if $\ct_t = 1-c^*$, then $\vy_t = v_{b^*,c^*}(x_t,\ct_t,\cf_t) = v_{b^*,c^*}(x_t,1-c^*,\hy_t) = \hy_t$. Therefore, when the optimal \dt strategy described above is used, if in the first round \dt selects $(\hy,\ct) = (1-b^*,1-c^*)$, it incurs no loss and does not change its strategy for the rest of the run, never trying $b^*, c^*$, and incurring a prediction mistake in every round. The probability that \dt tries $(\hy,\ct) = (1-b^*,1-c^*)$ in the first round is $1/4$. Hence, using these two debaters leads to $\ex[\Mi(T)] \geq T/4$.

Second, consider the optimal \df according to regret minimization. Consider the following strategy: In every round, \df selects $\cf_t = 1-\hy_t$. Then, if $\ct_t = 1-c^*$, we have $\vy_t = v_{b^*,c^*}(x_t,\ct_t,\cf_t) = v_{b^*,c^*}(x_t,1-c^*,1-\hy_t) \neq \hy_t$ and \df does not incur a loss. If  $\ct_t = c^*$, then this is an irrefutable round, hence \df may incur a loss but it does not accumulate regret. Thus, this strategy obtains zero regret and so it is optimal with respect to regret minimization. 

Using this strategy against the optimal \dt strategy described above leads to \dt incurring a loss with every pair except for $(\hy,\ct) = (b^*,c^*)$. Therefore, the expected number of rounds until it tries $(\hy,\ct) = (b^*,c^*)$ is $1.5$, leading to $\Mi(T) \leq 1.5$. This can also be derived from Claim~\ref{claim:regret}.
\end{proof}

 The results above motivate setting the objective of oppositional debaters to regret instead of loss. Nonetheless, as we show next, even with this objective, oppositional debaters can still fail where cooperative debaters succeed.



\subsection{Gap between cooperative and oppositional debaters}

The following example shows that oppositional debaters optimizing worst-case regret can fail to obtain a low number of prediction mistakes, even when cooperative debaters succeed. 


\begin{example}\label{ex:reggap}
Let $\X=\{e\}\cup\naturals$, where $e\notin\naturals$, and let
$\C=\{0,1\}$. Consider length-1 debate. For $b\in\{0,1\}$, define
\[
h_b(x)=
\begin{cases}
b,&x=e,\\
0,&x\in\naturals.
\end{cases}
\]
For each $b\in\{0,1\}$ and $s:\naturals\to\{0,1\}$, define
\[
v_{b,s}(x,p,r)=
\begin{cases}
b,&x=e,\\
0,&x\in\naturals\text{ and }(p=s(x)\text{ or }r=0),\\
1,&\text{otherwise}.
\end{cases}
\]
Set $\V=\{v_{b,s}:b\in\{0,1\},\ s:\naturals\to\{0,1\}\}$.
The following table illustrates the verifier output for $x \in \naturals$.
\[
\begin{array}{c|cc}
v_{b,s}(x,\ct,\cf) &\cf=0&\cf=1\\\hline
\ct=s(x)&0&0\\
\ct \ne s(x)&0&1
\end{array}
\]
%The certificate $p=s(x)$ forces label $0$, and $r=0$ prevents any
%certificate from forcing label $1$. At $e$, every entry is $b$.
Note that $v_{b,s}$ validate $h_b$.
\end{example}

\begin{theorem}\label{thm:reggap}
For Example~\ref{ex:reggap}, worst-case regret minimizer strategies for \df and \dt obtain $\tRe(T) \leq T/2$ and $\fRe(T) = 0$, and in the worst case, $\Mi(T) = T/2$.
On the other hand, $\dldim(\V) = 1$, hence deterministic cooperative debaters incur at most one prediction mistake. 
\end{theorem}












\iffalse
\subsection{(\sivan{work in progress) dimensions for oppositional debaters}}

To evaluate the worst-case regret of an oppositional debater against an adversary, we can define appropriate dimensions for \dt and \df.
For simplicity, we define this for debates of length $n=1$. \sivan{to add how this works for other $n$}
In this construction, similarly to above, edges correspond to all possible actions of the debater, and nodes correspond to worst-case actions of the environment and the other debater. Unlike the $\dldim$, in this case defining the dimension for each debater is straightforward, since their loss is directly observed. 

In a tree for \dt with verifier class $\V$, each tree node is of the form $(\cf,\vy,x)$, and each edge is labeled with $(\vy, \ct)$. The root node is $(\sqcup,\sqcup,x)$. Each node has exactly one down-directed edge for each possible \dt action $(\hy, \ct)$. Each path from root to leaf is consistent with some $v \in \V$. Moreover, each edge labeled $(\hy, \ct)$ ends in a node with $\vy \neq \hy$. The depth of such a tree is defined as the minimum over the length over all its paths. The dimension $\ctd(\V)$ is defined as the maximum depth of a tree of this type. \sivan{this could probably be more formal...}

\begin{claim}
    In the realizable setting, for any \sivan{deterministic} choices of \dt with verifier class $\V$, the environment can enforce \dt to incur a loss of $\ctd(\V)$ on \dt. Moreover, there exists a deterministic strategy for \dt that never incurs a loss of more than $\ctd(\V)$.
\end{claim}
In this case, the SOA construction is similar to that of \cite{oz2026discriminative}: In each round with domain instance $x$, let the current version space be $\V'$. Let 
$\V'(\hy,\ct,\cf) = \{v \in \V' \mid v(x,\ct,\cf) \neq \hy\}$. 
The action $(\hy,\ct)$ selected by $\dt$ is some action in   
$\argmin_{(\hy,\ct)} \max_{\cf}(\ctd(\V'(\hy,\ct,\cf)))$. It is easy to see that if $\ctd(\V') \geq 1$, then $\min_{(\hy,\ct)} \max_{\cf}(\ctd(\V'(\hy,\ct,\cf))) < \ctd(\V')$, otherwise one could construct a \dt tree for $\V'$ with a larger depth than $\ctd(\V')$. Thus, by induction, after at most $\ctd(\V)$ steps, $\ctd(\V') = 0$, indicating that \dt can avoid additional loss rounds by choosing its actions wisely. For the converse, it is easy to see that given a verifier class with a finite $\ctd(\V)$, an adversary can enforce $\ctd(\V)$ on any deterministic \dt by taking the actions in the tree nodes corresponding to the path chosen by \dt.


For \df and verifier class $\V$, we again construct a corresponding tree. %However, in this case, we are interested in the regret $\fRe$ and not the loss. This is because an adversary that sets $x,\ct$ \textbf{\sivan{xXXXXXXXXXXXXXXxxx}}
Each tree node is of the form $(\vy,x,\hy, \ct)$, representing the verifier label $\vy$ for the previous round, the new domain instance $x$ and the response $(\hy, \ct)$ of \dt to $x$. Each edge is labeled with some $\cf \in \C$, representing the action of \df. The root node is of the form $(\sqcup,x,\hy, \ct)$. Each node has exactly one down-directed edge for each possible \df action $\cf \in \C$. Each path from root to leaf is consistent with some $v \in \V$. 
Moreover, each edge starting at a node $(\vy,x,\hy, \ct)$ ends in a node with $\vy = \hy$, incurring a loss of $1$ for \df. The depth of such a tree is defined as the minimum over the length over all its paths. 


\sivan{TODO: define the depth with respect to the loss only, since this is what this debater optimizes. The dimension $\cfd(\V)$ is defined as the maximum depth of a tree of this type. }
\fi

\section{Non-Realizability with Cooperative Debaters}

In the $k$-non-realizable setting for $k \in \naturals$, we consider the same debate protocol as described above, except that in at most $k$ rounds, the environment may provide $\vy \neq v^*(x_t, \sc_t)$. In the cooperative setting, we show upper and lower bounds for learning in the $k$-non-realizable setting. Our analysis shows that debate learning is more adversely affected by feedback noise or model misspecification than standard online learning.

\paragraph{Separation between non-realizable and realizable learning}
\label{sec:agnostic-example}
First, we show a simple verifier class that is easy to learn in the realizable cooperative case (with randomized debaters), but not when there are misspecifications.

\begin{example}\label{example:sep}
    Consider debates of length one. Let $\X=\naturals$, $\C=\{0,1\}$, and let $\ct_t$ and $\cf_t$ be the certificates at step $t$. For $b\in\{0,1\}$, let $h_b$ be the constant hypothesis $h_b(x)=b$, and let $\H=\{h_0,h_1\}$. For every bit $b$ and every function $e:\X\to \C\times\C$, define
\[
v_{b,e}(x,\ct,\cf)=
\begin{cases}
1-b, & (\ct,\cf)=e(x),\\
b, & \text{otherwise},
\end{cases}
\]
To illustrate, suppose that $e(x) = (c_0,c_1)$, then the value of $v_{b,e}(x,\ct,\cf)$ is as follows.
\[
\begin{array}{c|cc}
v_{b,e}(x,\ct,\cf) & \cf=c_1 & \cf=1-c_1\\
\hline
\ct=c_0 & 1-b & b\\
\ct=1-c_0 & b & b
\end{array}
\]
Let $\V=\{v_{b,e}:b\in\{0,1\},\ e:\X\to \C\times\C \}$. 
%For each $x$, think of $v(x,\cdot,\cdot)$ as a $2\times2$ table with rows indexed by the \dt's certificate $\ct$ and columns by the \df's certificate $\cf$. The table of $v_{b,e}$ at $x$ has every entry equal to $b$, except the entry $e(x)$, which equals $1-b$. If $e(x)=\bot$, there is no such entry, and the whole table is $b$. 

Clearly, for any $b,e$, $v_{b,e}$ validates $h_b$, using the $p$ that gives $b$ for any $r$. 
\end{example}


\begin{theorem}\label{thm:examplesep}
For Example~\ref{example:sep}, the following two properties hold. In the realizable case with cooperative debaters, there is a randomized strategy that obtains $\ex[\Mi(T)]<9$ for every $T$, every $v^*\in\V$ and every sequence of input instances. In the unrealizable case, for every $k\ge1$ and every $T\ge 3k$, the environment 
can choose the inputs and draw $v^*\in\V$ and the misspecified rounds at random in advance, so that the expected number of misspecified rounds is $k$ and every strategy of the debaters has $\ex[\Mi(T)]\ge\frac32k$.
\end{theorem}

% The idea is that a random entry of a table in $\V$ provides the true label $b$ with probability $3/4$, while a misspecified round can provide $1-b$ regardless of the choice of $\ct,\cf$. If each round is misspecified with probability $1/3$, 
% the probability of obtaining the true $b$ in each round is $1/2$. Consider a sequence of non-repeating instances $x_1,x_2,\ldots x_T = 1,2,\ldots, T$. With the above misspecification strategy, it is impossible to gain any information on $b$ from the interaction. 


%  Let $x_t=t$ and $m=3k$, and consider the following random instance, drawn before the interaction. Sample $b$ uniformly and set $h^*=h_b$. Independently for each $t\le m$, with probability $1/3$ make round $t$ misspecified and set every entry of $v^*(x_t,\cdot,\cdot)$ to $b$, so that $\vy_t=1-b$ whatever the certificates. Otherwise, set every entry of $v^*(x_t,\cdot,\cdot)$ to $b$ except one uniformly random entry $e_t\in\C\times\C$, which we set to $1-b$. No later round is misspecified, and for every input $x>m$, every entry of $v^*(x,\cdot,\cdot)$ is $b$. Every table of $v^*$ has at most one entry that differs from $b$, so $v^*\in\V$ validates $h^*$. Let $K$ be the number of misspecified rounds, so $K\sim\mathrm{Binomial}(m,1/3)$. In particular, $\ex[K]=m/3=k$.

% Fix a strategy (potentially randomized) of the debaters, which may depend on $k$ and $T$. We show that $\ex[\Mi(m)]=m/2$. Let $I_{t-1}$ consist of the debaters' random bits and the transcript of the first $t-1$ rounds. On a misspecified round $\vy_t=1-b$, and on any other round $\vy_t=b$ unless $(\ct_t,\cf_t)=e_t$. Hence, for each $a\in\{0,1\}$ and any certificates, $\Pr(\vy_t=a\mid h^*=h_a,I_{t-1})=\frac23\cdot\frac34+\frac13\cdot0=\frac12$. So whichever of $h_0$ and $h_1$ is the true hypothesis, each $\vy_t$ is a fair coin independent of $I_{t-1}$, which means $\Pr(\hy_t\ne\ty_t)=1/2$.  Summing over $t\le m$ gives $\ex[\Mi(m)]=m/2$. Since $m=3k\le T$, we get $\ex[\Mi(T)]\ge\ex[\Mi(m)]=\frac32k$, while the expected number of misspecified rounds is $\ex[K]=k$.





\paragraph{A lower bound for general verifier classes}
Our next result shows that %for deterministic learners, 
misspecification can result in a loss that is linear in $k\cdot \dldim(\V)$.
\begin{theorem}\label{thm:kdlower}
    For every $d$ and every $n$, there exists a class $\V_{d,n}$ with $\dldim(\V_d)=d$ and such that every %deterministic 
    learner makes $\Omega(kd)$ mistakes in the $k$-non-realizable case.
\end{theorem}

\paragraph{An upper bound for the $k$-non-realizable setting}
We provide a matching mistake upper bound of $O(k\cdot \dldim(\V))$ for any verifier class $\V$, using a randomized strategy. \cite{oz2026discriminative} provides a general deterministic strategy for obtaining such an upper bound for interactive learning protocols. However, that strategy relies on the learner having access to its incurred loss, which is not available in debate learning. 
Our upper bound employs instead a simple randomized algorithm based on majority voting, whose structure is inspired by the Private Online Procedure of \cite{kaplan2023black}. 

%We note that this upper bound uses a randomized strategy, while the lower bound above holds for deterministic strategies. 
\begin{theorem}\label{thm:upperkd}
    Assume the $k$-non-realizable setting. For any verifier class $\V$ with $\dldim(\V) = d \in \naturals$, there is a randomized learning strategy for cooperative debaters that incurs at most $16kd$ mistakes in expectation. 
\end{theorem}


\section{Conclusion}

% \todo{}
In this work, we initiated the formal study of learning through debate. While debate has been used in training pipelines, no formal framework for understanding its effectiveness for \emph{learning} had been developed so far. Our analysis highlights some important considerations from a safety perspective. We show that the practice of setting the debaters up as opponents, rather than cooperating agents, can be sub-optimal. 
This finding suggests that developing training pipelines with cooperating, rather than competing agents, might be a promising avenue, which is yet underexplored.
Moreover, our results demonstrate the brittleness of learning through debate to noise in the verifier feedback or slight model misspecifications. This result shows that employing debating agents as facilitators for learning demands caution: if the verifier that these agents interact with is not fully reliable, the learning process may suffer disproportionately in accuracy.

\section*{AI Disclosure}

For this work, we used ChatGPT and Claude for brainstorming some proof ideas and some examples. All derivations and proofs in the paper were written and verified by the authors.


\iffalse
\section{Unrealizable oppositional setting \todo{complete this}}

In this case, in the oppositional setting we extend the regret to account for misspecified rounds. For \dt, we define $\tRe(T) = \tLo(T) - k$, where $k$ is the minimal value such that the sequence is $k$-non-realizable. For \df, we define 
$\fRe(T) = \fLo(T) - k - k'$, where $k$ is the minimal value such that the sequence is $k$-non-realizable and $k'$ is the number of realizable rounds that are irrefutable. 

\vinayak{For all three of $\Mi(T), \tLo(T)$, and $\fLo(T)$, we define the corresponding notions of regret, where we subtract from the terms, the smallest loss achievable with the knowledge of the true $v^*$. Since the knowledge of the true $v^*$ would give us zero mistakes, we define \emph{debate regret} as simply $\Re(T) = \Mi(T)$. The \dt can't know what rounds are misspecified even if it knows the true $v^*$, therefore \dt's regret is $\Re_T(T) = \tLo(T) - k$. To define \df's regret, note that there are two kinds of rounds where it cannot achieve a loss of 0 even if it knows the true $v^*$: a) if the round is misspecified, and b) if the \dt has played optimally. We call a round where \dt has played optimally an \emph{irrefutable} round. Thus we define $\Re_F(T) = \fLo(T) - \#\text{rounds that are misspecified or irrefutable}$. }

and the \emph{debate regret}
quad\vinayak{\text{define regret in comparison to optimal debaters(?)}}
irrefutable rounds
\dt is evaluated through its \emph{loss} and its \emph{regret}
\begin{align*}
 \tLo(T) = \sum_{t = 1}^T \indct[\vy_t \neq \hy_t]\quad\text{and}\quad \Re_T(T) = \tLo(T) - \#\text{misspecified rounds}
\end{align*}

\vinayak{Things are a bit more complicated in the agnostic case because two different kinds of things make it impossible to achieve 0 loss: 1) if there is misspecification, i.e., no hypothesis in the class assigns the true label, and b) if its an anachievable round, i.e., \dt has picked the ``correct'' certificate and so \df cannot possibly win.}


\vinayak{
\begin{claim}
  $\Re(T) \leq \tRe(T) + \fRe(T) + 2k$.
\end{claim}
\begin{proof}
Let $u$ be the number of irrefutable rounds that are not misspecified. On an irrefutable round, the \dt's certificates force $v^*$ to output $\hy_t$ whatever the \df does. Since $v^*$ validates $h^*$, this means that $\hy_t=h^*(x_t)$, so there are no mistakes on these $u$ rounds. In every round, exactly one of the debaters incurs a loss, so
\[
\Mi(T) \leq T-u = \tLo(T)+\fLo(T)-u.
\]
By the definition of $\Re_T(T)$, we have $\tLo(T)=\Re_T(T)+k$. The rounds that are misspecified or irrefutable are the $k$ misspecified rounds together with the $u$ irrefutable rounds that are not misspecified, so $\fLo(T)-u=\Re_F(T)+k$. Substituting both gives $\Re(T)=\Mi(T)\leq\Re_T(T)+\Re_F(T)+2k$.
\end{proof}
}


\sivan{I Think the agnostic separation example could work for opposing ones but need to think about it ... if an instance appears for the first time in this sequence, \dt proposes as a label the majority of the earlier feedback $\vy_1,\ldots,\vy_{t-1}$, and chooses its certificate/argument uniformly and independently. Otherwise, XXXXXXXXXXXXxx Since each table of $v^*$ has at most one entry equal to $1-b$,
the probability that the response is $1-b$ is either $1/2$ }

\fi

% \section*{Acknowledgments}
% Add acknowledgments here.

\bibliography{references}

\clearpage
\appendix

% \crefalias{section}{appendix} % uncomment if you are using cleveref



\section{Deferred proofs}

\begin{proof}[Proof of Lemma~\ref{lemma:cooperative_deterministic_upperbound}]
For an instance $x$, define the set of trees $\Delta_n(\V,x)$ to be the set of trees in $\Delta_n(\V)$ that have $x$ at the root.
Define the set of trees $\Delta_n(\V,x,\hy,\sc)$ to be the set of trees that result from taking any tree from $\Delta_n(\V,x)$ and removing all the edges (and corresponding sub-trees) coming out of $x$ except the one labeled $\hy,\sc$. 

First, we show that $\min_{\hy,\sc}Q_t(\hy,\sc)$ is non increasing in $t$. Consider a round $t$ in some run of SOA. We have
\begin{align*}
  &\min_{\hy,\sc}Q_t(\hy,\sc) = \min_{\hy,\sc} \max_{\vy\in\{0,1\}} ~\sup_{\T\in\Delta_n(\V_{(p_t,\hy, \sc, \vy)})} ~\min_{q\in\T^\rho} ~\max_{h\in \H_{(p_t,\hy, \sc, \vy)}\cap \H_q} (\mi(h, p_t) + \mi(h,q)+\indct[h(x_t) = \hy])\\
  &\leq
    \min_{\hy,\sc}  ~\sup_{\T\in\Delta_n(\V_{p_t},x_t,\hy,\sc)} ~\min_{q\in\T^\rho} ~\max_{h\in \H_{p_t}\cap \H_q} (\mi(h, p_t) + \mi(h,q))
\end{align*}
For a path $q\in\T^\rho$, denote $f(q) := \max_{h\in \H_{p_t}\cap \H_q} (\mi(h, p_t) + \mi(h,q))$, and let
\[
q_\T \in \argmin_{q\in\T^\rho} f(q).
\]
denote the path in tree $\T$ that minimizes the sum of mistakes an $h$ can incur in $q_\T$ and $p_t$ combined.
Note that if $\T \in \Delta_n(\V_{p_t},x_t,\hy,\sc)$ then $q_\T$ must start with $(x_t,\hy,\sc)$. Let $q'_\T$ be the rest of the path of $q_\T$ and write $q_\T = (x_t,\hy,\sc);q'_\T$. 
Then we have 
\begin{align*}
&\min_{\hy,\sc}Q_t(\hy,\sc)  \leq \min_{\hy,\sc}  ~\sup_{\T\in\Delta_n(\V_{p_t},x_t,\hy,\sc)} f(q_\T) = 
\min_{\hy,\sc}\sup_{\T\in\Delta_n(\V_{p_t},x_t,\hy,\sc)} f((x_t,\hy,\sc);q'_\T) =: R
%&= \min_{\hy,\sc}\sup_{\T\in\Delta_n(\V_{p_t},x_t)} f((x_t,\hy,\sc);q'_\T)
\end{align*}

Define $\T_{\hy,\sc} \in \argmax_{\T \in \Delta_n(\V_{p_t},x_t,\hy,\sc)}f((x_t,\hy,\sc);q'_\T)$ (if a maximum exists)\footnote{A similar argument can be made if there is no maximum, by using a sequence of $\T_{\hy,\sc}$ choices that are larger than a sequence converging to the supremum}
and let $\T^*$ be the tree constructed by joining each of the trees $\T_{\hy,\sc}$ at the root $x_t$. Note that $\T^* \in \Delta_n(\V_{p_t},x_t)$. Let $X_{\hy,\sc}$ be the set of paths that start with $(x_t, \hy,\sc)$.
Then
\[
\min_{q \in (\T^*)^\rho} f(q) = \min_{\hy,\sc} \min_{q \in \T^*\cap X_{\hy,\sc}} f(q) = \min_{\hy,\sc}f((x_t,\hy,\sc);q'_{\T_{\hy,\sc}}) = R.
\]
Hence,
\[
\sup_{\T \in \Delta_n(\V_{p_t},x_t)} \min_{q \in \T^\rho} f(q)  \geq \min_{q \in (\T^*)^\rho} f(q) \geq R \geq \min_{\hy,\sc}Q_t(\hy,\sc).
\]
We now upper bound the LHS.
\begin{align*}
&\sup_{\T \in \Delta_n(\V_{p_t},x_t)} \min_{q \in \T^\rho} f(q) = \sup_{\T\in\Delta_n(\V_{p_t},x_t)} ~\min_{q\in\T^\rho} ~\max_{h\in \H_{p_t}\cap \H_q} (\mi(h, p_{t}) + \mi(h,q))
  \\
  &\quad= \sup_{\T\in\Delta_n(\V_{(p_{t-1},\hy_{t-1}, \sc_{t-1}, \vy_{t-1})},x_t)} ~\min_{q\in\T^\rho} ~\max_{h\in \H_{(p_{t-1},\hy_{t-1}, \sc_{t-1}, \vy_{t-1})}\cap \H_q}
  \\
  &\qquad\qquad (\mi(h, p_{t-1}) + \mi(h,q)+\indct[h(x_{t-1}) = \hy_{t-1}])
  \\
  &\quad\leq \max_{\vy\in\{0,1\}} ~\sup_{\T\in\Delta_n(\V_{(p_{t-1},\hy_{t-1}, \sc_{t-1}, \vy)})} ~\min_{q\in\T^\rho} ~\max_{h\in \H_{(p_{t-1},\hy_{t-1}, \sc_{t-1}, \vy)}\cap \H_q}
  \\
  &\qquad\qquad (\mi(h, p_{t-1}) + \mi(h,q)+\indct[h(x_{t-1}) = \hy_{t-1}])
\\   
  &\quad=\min_{\hy,\sc}\max_{\vy\in\{0,1\}} ~\sup_{\T\in\Delta_n(\V_{(p_{t-1},\hy, \sc, \vy)})} ~\min_{q\in\T^\rho} ~\max_{h\in \H_{(p_{t-1},\hy, \sc, \vy)}\cap \H_q} (\mi(h, p_{t-1}) + \mi(h,q)+\indct[h(x_{t-1}) = \hy]) \\
  &\quad= \min_{\hy,\sc}Q_{t-1}(\hy,\sc).
\end{align*}
We conclude that $\min_{\hy,\sc}Q_{t}(\hy,\sc) \leq \min_{\hy,\sc}Q_{t-1}(\hy,\sc)$. Now consider $Q_1$. By definition, $\mi(h, p_1) = 0$ since $p_1$ includes no hypothesized labels, and 
similarly to the above,
\[
\min_{\hy,\sc}Q_1(\hy, \sc) = \max_{\vy\in\{0,1\}} ~\sup_{\T\in\Delta_n(\V_{(p_1,\hy, \sc, \vy)})} ~\min_{q\in\T^\rho} ~\max_{h\in \H_{(p_1,\hy, \sc, \vy)}\cap \H_q} (\mi(h,q) +\indct[h(x_t) \neq \hy]).
\]
Since $\V_{p_1} = \V$, we get
\begin{align*}
\min_{\hy,\sc}Q_1(\hy, \sc) &\leq \sup_{\T\in\Delta_n(\V_{p_1},x_1)} ~\min_{q\in\T^\rho} ~\max_{h\in \H_{p_1}\cap \H_q} (\mi(h, p_{1}) + \mi(h,q)) \\
&= \sup_{\T\in\Delta_n(\V,x_1)} ~\min_{q\in\T^\rho} ~\max_{h\in \H_q} (\mi(h,q)) \leq \dldim(\V, n).
\end{align*}
The latter follows from the definition of $\dldim(\V, n)$.
Lastly, to upper bound the total number of mistakes of the SOA run, note that for path $p_t$ taken by SOA, at any round $t$,
% \begin{align*}
% &\max_{h\in \H_{(p_t,\hy_t, \sc_t, \vy)}} (\mi(h, p_t)) \leq 
% \min_{\hy,\sc} \max_{\vy\in\{0,1\}} ~\sup_{\T\in\Delta_n(\V_{(p_t,\hy, \sc, \vy)})} ~\min_{q\in\T^\rho} ~\max_{h\in \H_{(p_t,\hy, \sc, \vy)}\cap \H_q} (\mi(h, p_t) + \mi(h,q)+\indct[h(x_t) = \hy]) \\
% &\quad= \min_{\hy,\sc}Q_{t}(\hy, \sc) \leq \dldim(\V, n).
% \end{align*}
% \ruth{I also don't follow that first inequality that you are explaining in the next sentence..don't think there is a mistake in the proof overall, but can't follow the formalism at this point.. Maybe the $\hy$ and $\sc$ in the very first term (before the first inequality) should be marked as the ones that the SOA actually chooses? and then after the first inequality the indices should go down to $t-1$? Maybe:
\begin{align*}
&\max_{h\in \H_{p_t}} (\mi(h, p_t)) \\
&\quad\leq
\min_{\hy,\sc} \max_{\vy\in\{0,1\}} ~\sup_{\T\in\Delta_n(\V_{(p_{t-1},\hy, \sc, \vy)})} ~\min_{q\in\T^\rho} ~\max_{h\in \H_{(p_{t-1},\hy, \sc, \vy)}\cap \H_q} (\mi(h, p_{t-1}) + \mi(h,q)+\indct[h(x_t) = \hy]) \\
&\quad= \min_{\hy,\sc}Q_{t-1}(\hy, \sc) \leq \dldim(\V, n).
\end{align*}


The first inequality follows by the realizability assumption, which implies that there is at least one $q$ with a non-empty $\H_q$. 
\end{proof}

\begin{proof}[Proof of Theorem~\ref{thm:reggap}]
Consider a cooperative learner: it can predict $\hy = 0$ for any $x_t \in \naturals$ without making any mistake. For $x_t = e$, it needs to make at most one mistake, since for $x_t = e$, $\vy = b$ is provided, revealing the label of $e$.

On the other hand, consider possible worst-case regret-optimizing debaters. 
First, define a regret-optimizing strategy for \df. This strategy refutes whenever possible: for any input $(x_t,\hy_t,\ct_t)$, it selects $\ct_t = 1-\hy_t$. Whenever $x \in \naturals$ and $\ct_t \neq s(x)$, this gives $\vy_t = v_{b,s}(x_t,\ct_t,\cf_t) = \cf_t \neq \hy_t$. Thus, its regret is zero and so optimal. 

Second, define a regret-optimizing strategy for \dt. For $x_t = e$, this strategy predicts a uniformly random label with $\ct_t = 0$. For $x_t \in \naturals$, it predicts $\hy = 0$ with a uniformly random $\ct_t$. Its expected loss on a given sequence is equal to half the number of rounds in which $e$ appears, thus $\tLo(T) = \tRe(T) \leq T/2$ and $T/2$ is also its worst-case regret over all sequences. This strategy is worst-case regret optimal, because any strategy obtains a regret of $T/2$ in the following scenario: If the environment selects $v^*$ uniformly at random, and presents distinct $x_t \in \naturals$ in each round, and if \df always selects $\ct_t = 1$, then \df always has a probability of $1/2$ to select $\hy_t = \indct[\ct_t\neq s(x_t)] = v^*(x_t, \ct_t, 1) = \vy_t$. Thus it suffers a loss and regret of $T/2$ in expectation. 

Now, observe that running these regret-optimizing \dt and \df against each other with input $e$ results in a prediction mistake. Hence, if $x_t = e$ in every round, $\Mi(T) = T/2$. 
\end{proof}

\begin{proof}[Proof of Theorem~\ref{thm:examplesep}]
We will consider debates of length one. Let $\X=\naturals$, $\C=\{0,1\}$, and let $\ct_t$ and $\cf_t$ be the certificates at step $t$. For $b\in\{0,1\}$, let $h_b$ be the constant hypothesis $h_b(x)=b$, and let $\H=\{h_0,h_1\}$. For every bit $b$ and every function $e:\X\to \C\times\C$, define
\[
v_{b,e}(x,\ct,\cf)=
\begin{cases}
1-b, & (\ct,\cf)=e(x),\\
b, & \text{otherwise},
\end{cases}
\]
To illustrate, suppose that $e(x) = (c_0,c_1)$, then the value of $v_{b,e}(x,\ct,\cf)$ is as follows.
\[
\begin{array}{c|cc}
v_{b,e}(x,\ct,\cf) & \cf=c_1 & \cf=1-c_1\\
\hline
\ct=c_0 & 1-b & b\\
\ct=1-c_0 & b & b
\end{array}
\]
Let $\V=\{v_{b,e}:b\in\{0,1\},\ e:\X\to \C\times\C$\}. 
%For each $x$, think of $v(x,\cdot,\cdot)$ as a $2\times2$ table with rows indexed by the \dt's certificate $\ct$ and columns by the \df's certificate $\cf$. The table of $v_{b,e}$ at $x$ has every entry equal to $b$, except the entry $e(x)$, which equals $1-b$. If $e(x)=\bot$, there is no such entry, and the whole table is $b$. 

Clearly, for any $b,e$, $v_{b,e}$ validates $h_b$, using the $p$ that gives $b$ for any $r$. %The exceptional entry, if there is one, lies in a single row, so the other row is all $b$. No row is all $1-b$, because that would need two entries that differ from $b$. So the \dt has a certificate that forces the label $b$ and none that forces $1-b$. Hence $h_{v_{b,e}}=h_b$, and $\H_\V=\H$.
Next, we give a randomized strategy for learning $h^*$ in the realizable case. Fix $v^*\in\V$ and a sequence of inputs, and let $h^*=h_{v^*}=h_b$. 


In each round, both \dt and \df choose their arguments $\ct,\cf$ uniformly and independently. \dt proposes as a label the majority of the earlier feedback $\vy_1,\ldots,\vy_{t-1}$, breaking ties arbitrarily. 
The values $\vy_t$ are independent, and each value equals $b$ with probability $3/4$. After $s$ rounds, the majority is wrong only if at most $s/2$ of the $s$ values equal $b$, while their expected number is $3s/4$. By Hoeffding's inequality, this has probability at most $e^{-2(s/4)^2/s}=e^{-s/8}$. Summing over rounds gives $\ex[\Mi(T)]\le\sum_{s\ge0}e^{-s/8}=1/(1-e^{-1/8})<9$. 

%\sivan{move this? The randomization is needed here. If the debaters' choices were deterministic, the environment could put the exceptional entry of each new table where they look.}

Next, we prove a lower bound in the unrealizable case. Suppose that $v^*$ is selected such that for each instance, $1-b$ appears at a uniformly random entry of the response matrix. Consider a sequence of non-repeating instances $x_1,x_2,\ldots x_T = 1,2,\ldots, T$, and assume that each round is misspecified with probability $1/3$, and that in such a round the response is always $1-b$ for any $\ct,\cf$. In rounds that are not misspecified, the probability of observing $b$ is $3/4$, while it is zero in misspecified rounds. Hence, the probability of observing $b$ is exactly $1/2$. It is thus impossible to gain any information on $b$ from the interaction, which implies that any prediction strategy of $\ct$ will result in an expected mistake rate of $1/2$. Since the rate of misspecifications is $1/3$, the expected number of misspecification rounds is $k = T/3$. Thus, the expected mistake rate is $\frac{3}{2}k$. 
\end{proof}


\iffalse

\begin{lemma}\label{lem:equilibrium}
    Let $\pi_\ct$ be the strategy of \dt out of all potential strategies $\Pi_\ct$ and let $\pi_\cf$ be the strategy of \df out of all potential strategies $\Pi_\cf$. Assume that the debate protocol runs for $T$ steps. Let $a$ be the strategy of the environment out of all possible strategies $A$. 
    Let $\tLo(a, \pi,\pi'), \fLo(a,\pi,\pi')$ be the losses of $\dt,\df$ respectively when the environment runs $a$, \dt runs $\pi$ and \df runs $\pi'$. 
    Suppose that 
    \begin{align*}
            &\pi_\ct \in \argmin_{\pi \in \Pi_\ct} \max_{a \in A} \tLo(a,\pi,\pi_\cf)\text{     and }
    \pi_\cf \in \argmin_{\pi \in \Pi_\cf} \max_{a \in A} \fLo(a,\pi_\ct,\pi).
    \end{align*}
    Then 
    \[
   \tLo(\pi_\ct, \pi_\cf)
    = \min_{\pi \in \Pi_\ct}\max_{\pi'\in \Pi_\cf}\tLo(\pi,\pi')
\quad \text{ and }\quad
\fLo(\pi_\ct, \pi_\cf)
 =    \min_{\pi' \in \Pi_\cf}\max_{\pi\in \Pi_\ct}\fLo(\pi,\pi').
    \]
    It follows that debaters that minimize their loss with respect to a worst-case opposing debater and worst-case environment are worst-case optimal also against a fixed opposing debater that minimizes its loss in the same way. 
\end{lemma}
\begin{proof}
    This claim follows because $\tLo = T-\fLo$, by a direct application of the saddle-point 
sufficiency condition \cite[see, e.g.,][]{boyd2004convex}. Let $\tLo(\pi,\pi') = \max_{a \in A} \tLo(a,\pi,\pi')$ and $\fLo(\pi,\pi') = \max_{a \in A} \fLo(a,\pi,\pi')$. For any $\pi \in \Pi_\ct,\pi' \in \Pi_\cf$, we have by the assumption that 

\[
\tLo(\pi_\ct, \pi') =
    T - \fLo(\pi_\ct, \pi')
    \leq 
    T - \fLo(\pi_\ct, \pi_\cf)
    \leq \tLo(\pi_\ct, \pi_\cf)
    \leq \tLo(\pi, \pi_\cf).
\]
Hence, by the saddle-point 
sufficiency condition,
\[
   \tLo(\pi_\ct, \pi_\cf)
    = \min_{\pi \in \Pi_\ct}\max_{\pi'\in \Pi_\cf}\tLo(\pi,\pi')
\]
Symmetrically,
\[
\fLo(\pi_\ct, \pi_\cf)
 =    \min_{\pi' \in \Pi_\cf}\max_{\pi\in \Pi_\ct}\fLo(\pi,\pi').
    \]
\end{proof}


\begin{proof}[Proof of Theorem~\ref{thm:kdlower}]
    Consider an infinite $\X$ and $\C = \{0,1\}^*$, the set of finite binary sequences.\tosca{still needs to be made finite because otherwise dimension not defined (for finite certificate space)}  Consider $\H = \{h_0,h_1\}$, where $h_a(x) = a$ for all $x\in \mathcal{X}, a\in\{0,1\}$. 
    We now will construct a verifier class, that is easy to learn in the realizable setting but hard to learn in the $k$-non-realizable setting.
  To ensure that the class we construct indeed contains only valid  verifiers, and supports $\mathcal{H}$ we will consider verifiers $v$, such that there exists $h\in \mathcal{H}$ and $b\in \C$ with:
  \begin{itemize}
      \item For all $x \in \mathcal{X}$, for $c_1 =1$, $\forall c_2 \dots c_{2n-2}\in \C, \dots, c_{2n-1}= b, \forall c_{2n}  $: $v(x, c_1, c_2, c_3,c_4, \dots, c_{2n-1}, c_{2n}) = h(x) $.
      \item There is $c_{\mathrm{out}},c_{\mathrm{lier}}, $ such that all $x \in \mathcal{X}$, for $c_1 =1$, $\forall c_2, \dots, c_{2n-2} \in C,  c_{2n-1} \in \C \setminus \{b\}$: $v(x, c_1. c_2, c_3,c_4, \dots, c_{2n-1}, c_{\mathrm{out}}) \neq v(x, c_1. c_2, c_3,c_4, \dots, c_{2n-1}, c_{\mathrm{lier}}) $% 1-h(x) $ and $c_1 =1$, $\forall c_2 \in \C, c_3 = 1, \forall c_4 \in \C , \dots, c_{2n-1}= 1- b, \forall c_{2n}   =c_{\mathrm{outlier}} $, $v(x, c_1. c_2, c_3,c_4, \dots, c_{2n-1}, c_{2n}) = h(x) $.
      \item For every $c_1 \in \C \setminus \{1\}$, for $c_2 = 0 $, for all $c_3, \dots, c_{2n} \in C$, we have $v(x,c_1, c_2, c_3, \dots, c_{2n}) = 0$; and for every $c_1 \in \C \setminus \{1\}$, for $c_2 = 1 $, for all $c_3, \dots, c_{2n} \in C$, we have $v(x,c_1, c_2, c_3, \dots, c_{2n}) = 1$. 
      \item For every $c_1\in \C \setminus \{0,1\}$ there is $l \in \{0,1\}$, such that for every $x\in\X$ and for every $c_2\in \naturals \setminus \{0,1\}$ and every $c_3, \dots, c_{2n} \in \naturals$: $v(x,c_1, \dots, c_{2n})= l $.
  \end{itemize}
    The first condition makes sure that there is always one successful debate strategy (in which the $\dt$ always chooses from $1$ in the first $n-1$ rounds and then chooses some certificate $b$ that is specific to $v$) which reveals the true label on every instance. 
    However, note that it is still hard for $\dt$ to learn that strategy, as all candidates for $b$ would have to be tried out (and there may not be feedback that distinguishes between getting the correct $b$ or not). 
    % However, formally this still ensures that $\V$ is a valid verifier class for $\H$.
    
    The second and third condition above ensure that each verifier in $\V$ validates exactly one hypothesis in $\mathcal{H}$.
    
    The last condition just states that if both $c_1$ and $c_2$ are not $0$ or $1$, then every $c_1$ defines a specific label $\mathcal{l}$ such that the rest of the dialogue will always evaluate to that label $l$.
    We will now use this specific label in order to encode the hypothesis $h$, in the following way.
    The idea is to encode a binary path of certificates (that work for every $x$, such that the label at the $d$-th element of that path will reveal whether the ground truth hypothesis is $h_0$ or $h_1$.
    % More formally, consider the binary representation of natural numbers, i.e. $\overline{100}=4$. 
    Now the class $\V$ will only contain verifiers $v$ for which there exists a sequence $a_1, \dots a_d \in \{0,1\}$
    \begin{itemize}
        \item For every $x\in \X$, for $c_1 = \dots =c_{2n} = 100$: $v(x, c_1, \dots, c_{2n}) = a_1$
        \item For every $x\in \X$, for $c_1 = \dots =c_{2n} = a_1100$, $v(x, c_1, \dots, c_{2n}) = a_2$
        \item $\dots$
        \item For every $x \in \X$ for $c_1 = \dots = c_{2n} = a_{r-1} \dots a_1 1 00$, we have $v(x, c_1, \dots, c_{2n}) = a_r$
        \item $\dots$
        \item For every $x \in \X$ for $c_1 = \dots = c_{2n} = a_{d-2} \dots a_1 1 00$, we have $v(x, c_1, \dots, c_{2n}) = a_d$
        \item For every $x \in \X$ for $c_1 = \dots = c_{2n} = a_{d-1} a_{d-2} \dots a_1 1 00$, we have $v(x, c_1, \dots, c_{2n}) = h_v(x)$.
    \end{itemize}
\sivan{is there a typo here with the $100$ numbers in the list above?}
    \paragraph{$\V$ has a mistake bound of $d$ in the realizable case} This can be shown by presenting a learner with mistake bound $d$. A learner can independent of input use certificate debates along the "path" that encodes the ground truth hypothesis. That is, in the first ground, the learner will use debate $c_1^1 = \dots = c_{2n}^1 = {100}$ and an arbitrary label $\hat{y}_0$ and receive output $a_1 : = v^* (x_1, c_1^1, \dots, c_{2n}^1)$. Then, in any subsequent round $t \leq d-1$, the learner will predict an arbitrary label and use the debate with $c_{1}^t = \dots = c_{2n}^t  = {a_{t-1} a_{t-2} \dots a_1 1 0 0}$ and receive label $a_{t} : = v^*(x_t,c_1^t, \dots, c_{2n}^t )$. 
 Then, in round $t = d $ the learner will use the debate with $c_{1}^d = \dots = c_{2n}^d  = {a_{d-1} a_{d-2} \dots a_1 1 0 0}$ and receive label $\ty : = v^*(x_t,c_1^d, \dots, c_{2n}^d )$.
 In every round $t > d$, the learner will then predict $\hat{y}_t= \ty$. According to the construction $\ty = h(x)$ for all $x\in \X$. Thus, this learner will make at most $d$ mistakes in the realizable case (and therefore  $\dldim(\V) \leq d$.)
\newcommand{\fat}{\textrm{fat}}
\paragraph{k-non-realizable case} We will argue that for every learner and there exists a $k$-non-realizable sequence sequence of $T= k\cdot d$ rounds, in which the learner is forced to make a mistake in at least half of the rounds. 
For this consider the following adversary behaviour: 
\begin{itemize}
    \item In every round $t$, present arbitrary $x_t \in \X$ and every $\bm{c} \in \C^{2n}$: 
    \begin{itemize}
        \item If $c_2^t = 1$ and $c_1 \neq 1$, output $\tilde{y}_{t} = 1$
        \item If $c_2^t = 0$ and $c_1 \neq 1$, output $\tilde{y}_{t} = 0$
        \item Otherwise, output output $\tilde{y}_{t} = 1$.
    \end{itemize}
\end{itemize}
We will argue that independent of the learner's behavior, the resulting sequence is $k$-consistent with a verifier $v^* \in \V$ such that $\sum_{t=1}^T \indct[h_{v^*}(x) \neq \hat{y}_t] \geq T/2$. \tosca{define k-consistent as consistent other than in $k$ exception rounds}\sivan{note we don't use "exceptions" in this paper, it's called misspecified rounds or feedback}
It is sufficient to argue that the sequence is $k$-consistent with both $h_0$ and $h_1$.
First we note, that the choice of sequence is always consistent with the verifier class $\mathcal{V}$.
We now argue that the predicted sequence will always be consistent with $h_1$, since the only way to get verifier feedback that is sufficient for reducing the version space for $\mathcal{H}$ in \tosca{finite/small number of rounds? depends on allowed certificate space which one} is through figuring out the sequence $a_{d-1} \dots a_1$ and getting verifier feedback $v(x, \bar{c}) = h^* (x) $ with $\bar{c} = c_1c_2 \dots c_{2n}$ with $c_1 = \dots = c_{2n} = a_{d-1} \dots a_1 1 0 0$. Since the adversary would pick output $\tilde{y}_t =1$ for any such choices of certificates, the sequence will remain consistent with hypothesis $h_1$.
Lastly we argue that the sequence is $k$-consistent with $h_0$.
For this we will define for every sequence a set of $k$ exceptions such that the sequence outside of those exceptions is consistent with $\V$ and $h_0$.
For $i = 3 \to d+2$ define $S_i := \{t : \bar{c}^t = c_1c_2 \dots c_{2n} \text{ with } c_1 = \dots = c_{2n} \text{ and } c_1 \text{ has length }i\}$.

We note, that for any sequence $s$, in order to reduce the version space to $\mathcal{H}_{\mathcal{V}_s}$ to less than two elements, the sequence must contain at least one (truthful) element from every $S_i$.

Furthermore, $\sum_{i=3}^{d+2} |S_i| \leq T = k d$.

Thus $\min_{i \in \{3, \dots, d+2}|S_i| \leq k$.

Thus, if the adversary declares all rounds in $ \arg\min_{S_i: i \in \{3, \dots, d+2}|S_i| $ as exception rounds, the version space defined by the remaining non-exception rounds includes both $h_0$ and $h_1$.
Thus both of the functions are $k$-consistent with any sequence resulting from the defined adversary and any learner. This concludes the proof.


% In order to prove this, it is sufficient to  show that for every set of debates $\bm{c}^1: =(c_1^1, \dots c_{2n}^1), \dots, \bm{c^T} = (c_1^T, \dots c_{2n}^T)$ and every verifier $v^* \in \V$, there exists an alternate verifier $v' \in \V$, such for at most $k$ rounds $i_1 \dots, i_k$, we have $v'(x_1,\bm{c}^{i_1} ) \neq v^*(x_1,\bm{c}^{i_`} ), \dots, v'(x_{i_k},\bm{c}^{i_k} ) \neq v^*(x_k,\bm{c}^{i_K} )$ but for all $x\in \X$, $h_{v^*}(x) \neq h_{v'}(x) $. If this is the case, then any learner that succeeds on $v^*$ (in every round) will fail on $v'$ in the $k$-non-realizable case. 

%We will now show that indeed, for every set of debates $\bm{c}^1: =(c_1^1, \dots c_{2n}^1), \dots, \bm{c^T} = (c_1^T, \dots c_{2n}^T)$ and every verifier $v^* \in \V$, there exists an alternate verifier $v' \in \V$, such for at most $k$ rounds $i_1 \dots, i_k$, we have $v'(x_1,\bm{c}^{i_1} ) \neq v^*(x_1,\bm{c}^{i_`} ), \dots, v'(x_{i_k},\bm{c}^{i_k} ) \neq v^*(x_k,\bm{c}^{i_K} )$ but for all $x\in \X$, $h_{v^*}(x) \neq h_{v'}(x) $.
%First we note that 
 
\end{proof}



\fi

\begin{proof}[Proof of Theorem~\ref{thm:kdlower}]


Let $\mathcal{X}=\{x_0\}$ and $\C = [4d]$.
Define the following verifiers. Let $v_0: \X \times \C^2 \to \{0,1\}$ defined by  $v_0(x_0,c_1,c_2) = 0 $ for every $c_1, c_2 \in \C$.
Furthermore for every $j \in [4d]$, let $v_j: \X \times \C^2 \to \{0,1\}$ be defined by 
\[v_j (x_0, c_1, c_2) = \begin{cases}
    1 & \text{ if } c_1=j \\
    1 & \text{ if } c_2=j \\
        0 & \text{ otherwise. } \\
\end{cases}\]

Now let $\mathcal{V}_d= \{v_j : j \in [4d] \} \cup \{v_0\}$.

We note that $v_0$ is a verifier for the all-zero hypothesis $h_0: \X \to \{0,1\}$, $h_0(x_0) = 0$.
Furthermore, we note that every $v_j$ is a verifier for the all-one hypothesis $h_1: \X \to \{0,1\}$, $h_1(x_0) = 1$. We set $\H = \{h_0, h_1\}$.

\begin{claim}
 $\H$ is learnable through debate of length $1$ through access to $\mathcal{V}_d$ in the realizable case with a mistake bound of $d$.
\end{claim}
We will argue that the following learning strategy is successful: In round $t$, predict the label $\hat{y}_t= 1$ if $t$ is odd and $\hat{y}_t= 0$ if $t$ is even. As certificates $(c_1^t, c_2^t)$ choose $(2t-1,2t)$.
After $2d$ rounds, if the learner has seen one round $r$ with verifier feedback $\tilde{y}_r = 1$, keep predicting $\hat{y}_t = 1$ for $t > 2d$.
After $2d$ rounds, if the learner has only seen verifier feedback $\tilde{y}_1 = \dots = \tilde{y}_{2d} = 0$, keep predicting $\tilde{y}_t = 0$ for $t > 2d$.

We are in the realizable case, so we know the true verifier is either $v_j$ for some $j\in [4d]$ or it is $v_0$.

If the true verifier is $v^* = v_j$ for some $j\in [4d]$, then in one of the rounds from $1$ to $2d$, i.e., round $\lceil \frac{j}{2}\rceil$, the verifier feedback will have been $\tilde{y}_{\lceil\frac{j}{2}\rceil} =1$.
In this case the learner will have made $d$ mistakes in rounds $1$ to $2d$ and no subsequent mistakes as the learner kept predicting $1$ after $2d$ rounds and the true verifier $v_j$ verifies the hypothesis $h_1$.

If the true verifier is $v^* = v_0$, then the true hypothesis is $h_0$. Furthermore, in rounds from $1$ to $2d$ the verifier feedback received by the learner is $\vy_1 = \dots = \vy_{2d} = 0$.
In this case the learner will have made $d$ mistakes in rounds $1$ to $2d$ and no subsequent mistakes as the learner kept predicting $0$ after $2d$ which is consistent with the true hypothesis $h_0$.

Thus this learning strategy makes at most $d$ mistakes.

\begin{claim}
    Every randomized learner for $\H$ through access to $\V_d$ via debate of length $1$ makes at least $\Omega(kd)$ mistake in expectation. 
\end{claim}

To show this we will fix a randomized adversary that picks the true verifier in advance, and can decide to give a different response in up to $k$ rounds after having seen the learners action $(\hat{y}_t, c_1^t, c_2^t)$ for that round. 
The adversary picks $v^*$ as $v_0$ with probability $1/2$ and for every $j\in [4d]$ picks every $v_j$ with probability $1/8d$.
Thus the true hypothesis is $h^* = h_0$ with probability $1/2$ and $h^* = h_1$ with probability $1/2$.

Now for any sequence, the adversary will flip to $0$ the first $k$ verifier responses that would have been $1$ if it responded according to the true verifier. 


Now consider any fixed (potentially randomized) learning  strategy $\mathcal{A}: (\X \times \Y\times \C^2 \times \Y)^*\times \X \to \Y \times \C^2$, that in every round $t$ maps the history of previous examples
\[
H = ((x_1, \hy_1, \sc_1, \vy_1 ), \dots, (x_{t-1}, \hy_{t-1}, \sc_{t-1}, \vy_{t-1} ))
\]
and current example $x_t$ to the next action $(\hy_t, \bar{c}_t)$.

We will call a history $H_{t}$ with $\vy_1 = \dots  = \vy_{d-1} = 0$ a \emph{null-history}. We will analyse both the learners behaviour given a null-history and the likelihood of obtaining a null-history in round $t$. We will combine these observations in order to show that the likelihood of obtaining a null-history is high and that when given a null-history in most cases the learner will not be able to successfully distinguish between verifiers for hypothesis $h_0$ and verifiers for hypothesis $h_1$.

For every $j$, we consider the count $N_j(t) = |\{\bar{c}_r: 1\leq r \leq t-1 \text{ with } c_1^t=j \text{ or } c_2^t=j\}|$.
%with the defined adversary the probability of obtaining of a sequence resulting in a null history is high. 

We note that given a null-history, if $v^* = v_0$, then independent of the learner's output in round $t$, we have $\ty_t = 0$ with probability $1$.
Furthermore, given a null history if $v^* = v_j$ then the adversary will output
\[ \vy_t = \begin{cases}
    0 &\text{if } N_{j}(t) \leq k \\
        1 &\text{otherwise.}
\end{cases}\]
with probability $1$.
% \ruth{why can we make an assumptiopn on the learner's behavior?} \tosca{ should be fized. I wrote learner instead of adversary again. I'm sorry}\ruth{please use ``vy'' for verifier respons, thanks ;)}

Thus, the only case in which $H_t$ is a nullhistory but $H_{t+1}$ is not a null-history is the case where
$v^* = v_j$ and $ N_j(t) = k $.
It follows that $H_t$ is a nullhistory unless $v^* =v_j$ with $N_j(t) >k$.

%\sout{We note that $H_1$ is a nullhistory. Thus by induction the only case in which $H_t$ is not a nullhistory, is one in which $v^* = v_j$ and there is $1 \leq r \leq t$  with $N_j(r) > k $. Since $N_j(r) \leq N_t(r)$,} 
% It follows that $H_t$ is a nullhistory unless $v^* =v_j$ with $N_j(t) >k$.
% \ruth{Not fully following these two sentences. Given the previous argument, can we reduce it to just: ``It follows that $H_t$ is a nullhistory unless $v^* =v_j$ with $N_j(t) >k$.''} \tosca{If you think the argument is clear, it's fine. I wanted to note that we always start with a nullhistory. Which is why in order to at some point not have a nullhistory, we have to transition to a nullhistory} 
% \ruth{Just given the sentence right before, but i made a change there, that might be crucial, the change from inequality to equality, is that one correct?}
% \tosca{yes that looks correct. you are right that this actually describtes the transition better. Let's change it thw way you propose} \tosca{the thing you state is correct. I wonder if the argument should be elaborated. But I'll trust your judgement} \tosca{feel free to delete this dialogue}
%Let $P_{r,j, H_{r}} = \mathbb{P}_{\mathcal{A}(H_{r}, x_0) = (\hat{y}, c_1^r, c_2^r) \text{ with } c_1^r = j \text{ or } c_2^r = j}$.
%Note that for every $r$ and every $H_r$, $\sum_{j=1}^{4d} P_{r,j, H_{r}} = 1$.
%Thus, for every choice of $H_1, \dots, H_t$:
%\[ \sum_{r=1}^t\sum_{j=1}^{4d} P_{r,j, H_{r}} = t. \]


We note that for every sequence of $t$ rounds, we have
\[
\sum_{j=1}^{4d} N_j(t) \leq 2t-2,
\]
since in every round $r \in [t]$ the learner can only increase the count for two different $j_1,j_2$ , namely $j_1= c_1^r$ and $j_2 = c_2^r$. 
%\tosca{I realized it's only an upper bound anyway, so just changed it. Then we don't need thereminerd} \tosca{I mainly wanted the convey the point to you. Feel free to judge how it comes across most efficiently}
%\tosca{ elaborating:
%\[
%\sum_{j=1}^{4d} N_j(t) =  \sum_{j=1}^{4d} \sum_{r=1}^{t} (N_j(r) - N_j(r-1)) =  \sum_{r=1}^{t} \sum_{j=1}^{4d} (N_j(r) - N_j(r-1)) \leq   \sum_{r=1}^{t} 2  \leq 2d.
%\]
%} 
Thus,
\[
\sum_{j=1}^{4d} \indct[N_j(t) > k] \leq \lfloor \frac{2t}{k} \rfloor
\]

Thus, after $t$ rounds for every sequence of choices by the learner, there are at most $\frac{2t}{k}$ different possible verifiers from $\{v_j: j\in [4d]\}$ that if the adversary picked it as the true verifier, it resulted in interactions between adversary and learner results in a history $H_t$ that is not a null-history.

The probability that the adversary picked such $v_j$ is therefore at most
\[ \frac{\frac{2t}{k}}{8d} = \frac{t}{4kd}.\]


We note, that given a null-history $H_t$ every learner makes a mistake with probability $1/2$ in the next round. That is because if the precition is $0$, it will fail in the case in which $v^*= v_j$ for $j\in [4d]$, which happens with probability $1/2$. And if the learner predicts $1$, it will fail in the case in which $v^* =v_0$, which happens with probability $1/2$.
Thus for every round $t$, 
\[\mathbb{E} [\indct[\hy_r \neq h_v(x)]|H_r \text{ is nullhistory}]  \geq 1/2\]

Using the probability of the sequence being a non-nullhistory at time $t$ together with the above inequality for the prediction error, given that the history is a nullhistory, we obtain the following:

\begin{align*}
   &\mathbb{E}[\sum_{r=1}^t  \indct[\hy_r \neq h_v(x)]]\\
   &\geq
   \mathbb{E} [\sum_{r=1}^t  \indct[\hy_r \neq h_v(x)|H_r \text{ is nullhistory}] \indct[H_r \text{ is nullhistory}]] \\
 & \geq \sum_{r=1}^t \frac{1}{2} \cdot \left(1-\frac{r}{4kd} \right)
   \geq \frac{t}{2}-\frac{t(t+1)}{16kd}  = \frac{t(8dk -t-1)}{16kd}.
\end{align*}

For $t = 4kd $ the above evaluates to $kd -\frac{1}{4} \in \Omega(kd)$. This proves the lower bound.
\end{proof}





\iffalse

    Consider an infinite $\X$ and $\C = \{0,1\}^*$, the set of finite binary sequences.\tosca{still needs to be made finite because otherwise dimension not defined (for finite certificate space)}  Consider $\H = \{h_0,h_1\}$, where $h_a(x) = a$ for all $x\in \mathcal{X}, a\in\{0,1\}$. 
    We now will construct a verifier class, that is easy to learn in the realizable setting but hard to learn in the $k$-non-realizable setting.
  To ensure that the class we construct indeed contains only valid  verifiers, and supports $\mathcal{H}$ we will consider verifiers $v$, such that there exists $h\in \mathcal{H}$ and $b\in \C$ with:
  \begin{itemize}
      \item For all $x \in \mathcal{X}$, for $c_1 =1$, $\forall c_2 \dots c_{2n-2}\in \C, \dots, c_{2n-1}= b, \forall c_{2n}  $: $v(x, c_1, c_2, c_3,c_4, \dots, c_{2n-1}, c_{2n}) = h(x) $.
      \item There is $c_{\mathrm{out}},c_{\mathrm{lier}}, $ such that all $x \in \mathcal{X}$, for $c_1 =1$, $\forall c_2, \dots, c_{2n-2} \in C,  c_{2n-1} \in \C \setminus \{b\}$: $v(x, c_1. c_2, c_3,c_4, \dots, c_{2n-1}, c_{\mathrm{out}}) \neq v(x, c_1. c_2, c_3,c_4, \dots, c_{2n-1}, c_{\mathrm{lier}}) $% 1-h(x) $ and $c_1 =1$, $\forall c_2 \in \C, c_3 = 1, \forall c_4 \in \C , \dots, c_{2n-1}= 1- b, \forall c_{2n}   =c_{\mathrm{outlier}} $, $v(x, c_1. c_2, c_3,c_4, \dots, c_{2n-1}, c_{2n}) = h(x) $.
      \item For every $c_1 \in \C \setminus \{1\}$, for $c_2 = 0 $, for all $c_3, \dots, c_{2n} \in C$, we have $v(x,c_1, c_2, c_3, \dots, c_{2n}) = 0$; and for every $c_1 \in \C \setminus \{1\}$, for $c_2 = 1 $, for all $c_3, \dots, c_{2n} \in C$, we have $v(x,c_1, c_2, c_3, \dots, c_{2n}) = 1$. 
      \item For every $c_1\in \C \setminus \{0,1\}$ there is $l \in \{0,1\}$, such that for every $x\in\X$ and for every $c_2\in \naturals \setminus \{0,1\}$ and every $c_3, \dots, c_{2n} \in \naturals$: $v(x,c_1, \dots, c_{2n})= l $.
  \end{itemize}
    The first condition makes sure that there is always one successful debate strategy (in which the $\dt$ always chooses from $1$ in the first $n-1$ rounds and then chooses some certificate $b$ that is specific to $v$) which reveals the true label on every instance. 
    However, note that it is still hard for $\dt$ to learn that strategy, as all candidates for $b$ would have to be tried out (and there may not be feedback that distinguishes between getting the correct $b$ or not). 
    % However, formally this still ensures that $\V$ is a valid verifier class for $\H$.
    
    The second and third condition above ensure that each verifier in $\V$ validates exactly one hypothesis in $\mathcal{H}$.
    
    The last condition just states that if both $c_1$ and $c_2$ are not $0$ or $1$, then every $c_1$ defines a specific label $\mathcal{l}$ such that the rest of the dialogue will always evaluate to that label $l$.
    We will now use this specific label in order to encode the hypothesis $h$, in the following way.
    The idea is to encode a binary path of certificates (that work for every $x$, such that the label at the $d$-th element of that path will reveal whether the ground truth hypothesis is $h_0$ or $h_1$.
    % More formally, consider the binary representation of natural numbers, i.e. $\overline{100}=4$. 
    Now the class $\V$ will only contain verifiers $v$ for which there exists a sequence $a_1, \dots a_d \in \{0,1\}$
    \begin{itemize}
        \item For every $x\in \X$, for $c_1 = \dots =c_{2n} = 100$: $v(x, c_1, \dots, c_{2n}) = a_1$
        \item For every $x\in \X$, for $c_1 = \dots =c_{2n} = a_1100$, $v(x, c_1, \dots, c_{2n}) = a_2$
        \item $\dots$
        \item For every $x \in \X$ for $c_1 = \dots = c_{2n} = a_{r-1} \dots a_1 1 00$, we have $v(x, c_1, \dots, c_{2n}) = a_r$
        \item $\dots$
        \item For every $x \in \X$ for $c_1 = \dots = c_{2n} = a_{d-2} \dots a_1 1 00$, we have $v(x, c_1, \dots, c_{2n}) = a_d$
        \item For every $x \in \X$ for $c_1 = \dots = c_{2n} = a_{d-1} a_{d-2} \dots a_1 1 00$, we have $v(x, c_1, \dots, c_{2n}) = h_v(x)$.
    \end{itemize}
\sivan{is there a typo here with the $100$ numbers in the list above?}
    \paragraph{$\V$ has a mistake bound of $d$ in the realizable case} This can be shown by presenting a learner with mistake bound $d$. A learner can independent of input use certificate debates along the "path" that encodes the ground truth hypothesis. That is, in the first ground, the learner will use debate $c_1^1 = \dots = c_{2n}^1 = {100}$ and an arbitrary label $\hat{y}_0$ and receive output $a_1 : = v^* (x_1, c_1^1, \dots, c_{2n}^1)$. Then, in any subsequent round $t \leq d-1$, the learner will predict an arbitrary label and use the debate with $c_{1}^t = \dots = c_{2n}^t  = {a_{t-1} a_{t-2} \dots a_1 1 0 0}$ and receive label $a_{t} : = v^*(x_t,c_1^t, \dots, c_{2n}^t )$. 
 Then, in round $t = d $ the learner will use the debate with $c_{1}^d = \dots = c_{2n}^d  = {a_{d-1} a_{d-2} \dots a_1 1 0 0}$ and receive label $\ty : = v^*(x_t,c_1^d, \dots, c_{2n}^d )$.
 In every round $t > d$, the learner will then predict $\hat{y}_t= \ty$. According to the construction $\ty = h(x)$ for all $x\in \X$. Thus, this learner will make at most $d$ mistakes in the realizable case (and therefore  $\dldim(\V) \leq d$.)
\newcommand{\fat}{\textrm{fat}}
\paragraph{k-non-realizable case} We will argue that for every learner and there exists a $k$-non-realizable sequence sequence of $T= k\cdot d$ rounds, in which the learner is forced to make a mistake in at least half of the rounds. 
For this consider the following adversary behaviour: 
\begin{itemize}
    \item In every round $t$, present arbitrary $x_t \in \X$ and every $\bm{c} \in \C^{2n}$: 
    \begin{itemize}
        \item If $c_2^t = 1$ and $c_1 \neq 1$, output $\tilde{y}_{t} = 1$
        \item If $c_2^t = 0$ and $c_1 \neq 1$, output $\tilde{y}_{t} = 0$
        \item Otherwise, output output $\tilde{y}_{t} = 1$.
    \end{itemize}
\end{itemize}
We will argue that independent of the learner's behavior, the resulting sequence is $k$-consistent with a verifier $v^* \in \V$ such that $\sum_{t=1}^T \indct[h_{v^*}(x) \neq \hat{y}_t] \geq T/2$. \tosca{define k-consistent as consistent other than in $k$ exception rounds}\sivan{note we don't use "exceptions" in this paper, it's called misspecified rounds or feedback}
It is sufficient to argue that the sequence is $k$-consistent with both $h_0$ and $h_1$.
First we note, that the choice of sequence is always consistent with the verifier class $\mathcal{V}$.
We now argue that the predicted sequence will always be consistent with $h_1$, since the only way to get verifier feedback that is sufficient for reducing the version space for $\mathcal{H}$ in \tosca{finite/small number of rounds? depends on allowed certificate space which one} is through figuring out the sequence $a_{d-1} \dots a_1$ and getting verifier feedback $v(x, \bar{c}) = h^* (x) $ with $\bar{c} = c_1c_2 \dots c_{2n}$ with $c_1 = \dots = c_{2n} = a_{d-1} \dots a_1 1 0 0$. Since the adversary would pick output $\tilde{y}_t =1$ for any such choices of certificates, the sequence will remain consistent with hypothesis $h_1$.
Lastly we argue that the sequence is $k$-consistent with $h_0$.
For this we will define for every sequence a set of $k$ exceptions such that the sequence outside of those exceptions is consistent with $\V$ and $h_0$.
For $i = 3 \to d+2$ define $S_i := \{t : \bar{c}^t = c_1c_2 \dots c_{2n} \text{ with } c_1 = \dots = c_{2n} \text{ and } c_1 \text{ has length }i\}$.

We note, that for any sequence $s$, in order to reduce the version space to $\mathcal{H}_{\mathcal{V}_s}$ to less than two elements, the sequence must contain at least one (truthful) element from every $S_i$.

Furthermore, $\sum_{i=3}^{d+2} |S_i| \leq T = k d$.

Thus $\min_{i \in \{3, \dots, d+2}|S_i| \leq k$.

Thus, if the adversary declares all rounds in $ \arg\min_{S_i: i \in \{3, \dots, d+2}|S_i| $ as exception rounds, the version space defined by the remaining non-exception rounds includes both $h_0$ and $h_1$.
Thus both of the functions are $k$-consistent with any sequence resulting from the defined adversary and any learner. This concludes the proof.


% In order to prove this, it is sufficient to  show that for every set of debates $\bm{c}^1: =(c_1^1, \dots c_{2n}^1), \dots, \bm{c^T} = (c_1^T, \dots c_{2n}^T)$ and every verifier $v^* \in \V$, there exists an alternate verifier $v' \in \V$, such for at most $k$ rounds $i_1 \dots, i_k$, we have $v'(x_1,\bm{c}^{i_1} ) \neq v^*(x_1,\bm{c}^{i_`} ), \dots, v'(x_{i_k},\bm{c}^{i_k} ) \neq v^*(x_k,\bm{c}^{i_K} )$ but for all $x\in \X$, $h_{v^*}(x) \neq h_{v'}(x) $. If this is the case, then any learner that succeeds on $v^*$ (in every round) will fail on $v'$ in the $k$-non-realizable case. 

%We will now show that indeed, for every set of debates $\bm{c}^1: =(c_1^1, \dots c_{2n}^1), \dots, \bm{c^T} = (c_1^T, \dots c_{2n}^T)$ and every verifier $v^* \in \V$, there exists an alternate verifier $v' \in \V$, such for at most $k$ rounds $i_1 \dots, i_k$, we have $v'(x_1,\bm{c}^{i_1} ) \neq v^*(x_1,\bm{c}^{i_`} ), \dots, v'(x_{i_k},\bm{c}^{i_k} ) \neq v^*(x_k,\bm{c}^{i_K} )$ but for all $x\in \X$, $h_{v^*}(x) \neq h_{v'}(x) $.
%First we note that 
 
\end{proof}
\fi


\begin{proof}[Proof of Theorem~\ref{thm:upperkd}]
The strategy works as follows: 
\begin{itemize}
    \item Initialize $4k$ copies of SOA for $\V$.
    \item In each round $t$
    \begin{itemize}
        \item Set $\hy_t$ to the majority label proposed by the $4k$ copies.
        \item Set $\sc_t$ to the certificate $\sc$ suggested by one of the copies $i_t$, where $i_t$ selected uniformly at random from $[4k]$. Get $\vy_t$ and feed it to the same copy.
    \end{itemize}
\end{itemize}
Call a copy ``uncorrupted at time $t$'' if it has not been provided with a corrupted feedback until round $t$. Since there can be at most $k$ misspecified rounds, at least $3k$ copies are uncorrupted at each time $t$. For round $t$, let $M_t = \indct[\hy_t \neq y_t]$. Let $\hy_t^i$ be the label proposed by copy $i$ at round $t$, and define $Z_t = \indct[\hy_t^{i_t} \neq y_t \wedge i_t \text{ is uncorrupted at time }$t$]$. Any copy $i_t$ can make at most $d$ mistakes in rounds in which its copy is updated, as long as it is uncorrupted. Hence, $\sum_t Z_t \leq 4kd$. On the other hand, if $M_t = 1$, then at least $2k$ copies $i$ have $\hy_t^{i} \neq y_t$ in round $t$, and at least $k$ of these are uncorrupted at time $t$. 
Therefore, for any history $H_t$ of the run until round $t$, $\Pr(Z_t = 1 \mid M_t = 1, H_t) \geq 1/4$. 
Thus,
\[
\ex[Z_t \mid H_t] \geq \Pr(Z_t = 1 \mid M_t = 1, H_t)\cdot \Pr(M_t = 1 \mid H_t) \geq \Pr(M_t = 1 \mid H_t)/4 = \ex[M_t \mid H_t]/4.
\]
Taking expectation over $H_t$, we get $\ex[Z_t] \geq \ex[M_t = 1]/4$.
Therefore,
\[
4kd \geq \ex[\sum_t Z_t] = \ex[\sum_t M_t]/4.
\]
It follows that $\ex[\sum_t M_t] \leq 16kd$.
\end{proof}




\end{document}





\section{\textcolor{red}{Some things we want to do}}

\begin{itemize}
    \item kd lower bound  - Tosca
    \item Lemma 6 - show that infinite certificate class doesn't work - Tosca?
    \item oppositional case dimensions - Sivan, working on it, regret vs loss issue
    \item oppositional case agnostic - ??
    \item is online learning always easier? - Tosca?
    \item regret minimizing which is locally optimal would kill the gap example, does it work more generally? i.e. is there an example where all collaborative solutions are not regret minimizing?
\end{itemize}
